% !TeX program = xelatex
% 由扫描件人工逐行转录。每个 \BookLine 对应原书一行，可直接编辑。
\documentclass[10pt]{article}
\usepackage[paperwidth=140mm,paperheight=210mm,margin=16mm]{geometry}
\usepackage{fontspec,xeCJK,graphicx,xcolor,tikz}
\usepackage[absolute,overlay]{textpos}
\usepackage[unicode,pdfpagelayout=TwoPageRight,pdftitle={阿尔刻提斯之歌（扫描部分重排）},pdfauthor={}]{hyperref}
\setmainfont{Liberation Serif}
\IfFileExists{/usr/local/share/fonts/f/FZShuSong_Z01_Regular.ttf}{
  \setCJKmainfont{FZShuSong_Z01_Regular.ttf}[Path=/usr/local/share/fonts/f/]
}{\setCJKmainfont{Noto Serif CJK SC}}
\IfFileExists{/usr/local/share/fonts/f/FZKai_Z03_Regular.ttf}{
  \newCJKfontfamily\kaifont{FZKai_Z03_Regular.ttf}[Path=/usr/local/share/fonts/f/]
}{\newCJKfontfamily\kaifont{FangSong}}
\IfFileExists{/usr/local/share/fonts/f/FZHei_B01_Regular.ttf}{
  \newCJKfontfamily\heifont{FZHei_B01_Regular.ttf}[Path=/usr/local/share/fonts/f/]
}{\newCJKfontfamily\heifont{Noto Sans CJK SC}}
\pagestyle{empty}
\setlength{\parindent}{0pt}
\setlength{\TPHorizModule}{1mm}
\setlength{\TPVertModule}{1mm}
\textblockorigin{0mm}{0mm}
\xeCJKsetup{PunctStyle=plain}
\newsavebox{\rowbox}
% y（毫米）, 缩进（毫米）, 字体, 是否铺满行, 文字。
\newcommand{\BookLine}[5]{%
  \begin{textblock*}{108mm}(16mm,#1mm)
  \fontsize{10.5}{14}\selectfont #3%
  \hspace*{#2mm}%
  \begingroup
  \ifnum#4=1 \renewcommand\CJKglue{\hskip 0pt plus 1fill}\fi
  \sbox{\rowbox}{#5}%
  \ifdim\wd\rowbox>\dimexpr108mm-#2mm\relax
    \resizebox{\dimexpr108mm-#2mm\relax}{!}{\usebox{\rowbox}}%
  \else
    \ifnum#4=1
      \makebox[\dimexpr108mm-#2mm\relax][s]{#5}%
    \else\usebox{\rowbox}\fi
  \fi
  \endgroup
  \end{textblock*}%
}
\newcommand{\Vertical}[5]{%
  \begin{textblock*}{8mm}(#1mm,#2mm)
    #3\fontsize{#4}{#5}\selectfont\centering
}
\newcommand{\EndVertical}{\end{textblock*}}
\newcommand{\BookNumber}[1]{%
  \begin{textblock*}{140mm}(0mm,192mm)
  \centering\fontsize{12}{12}\selectfont\itshape #1
  \end{textblock*}
}
\newcommand{\ActOrnament}{%
  \Vertical{66}{28}{}{14}{21}第\\一\\幕\EndVertical
  \begin{textblock*}{140mm}(0mm,63mm)
    \centering\textcolor{black!65}{\fontsize{10}{10}\selectfont\textbullet}
  \end{textblock*}
  \begin{textblock*}{75mm}(38mm,96mm)
  \begin{tikzpicture}[x=1mm,y=1mm,line width=.35pt,draw=black!65]
    % 后方拱门。
    \draw (16,0) -- (57,0) -- (57,39)
      .. controls (57,48) and (53,55) .. (48,59);
    \draw (16,0) -- (16,5);
    % 前方拱门。
    \draw (0,5) -- (48,5) -- (48,49)
      arc[start angle=0,end angle=180,radius=24] -- cycle;
    % 沿拱门上方逐字排出书名，不依赖扫描图。
    \foreach \letter [count=\i from 0] in {T,h,e,A,l,c,e,s,t,i,a,d}{%
      \pgfmathsetmacro{\a}{153-\i*11.45}
      \node[rotate=\a-90,text=black!80,font=\fontsize{9}{9}\selectfont]
        at ({24+35*cos(\a)},{49+35*sin(\a)}) {\letter};
    }
  \end{tikzpicture}
  \end{textblock*}
}
\begin{document}

% 原书页码：空白
\null
\newpage
% 原书页码：目录
\null
\Vertical{15}{28}{}{14}{21}目\\录\EndVertical
\BookLine{88.00}{29.52}{}{0}{前言 / 1}
\BookLine{105.00}{29.52}{}{0}{人物 / 20}
\BookLine{122.00}{29.52}{}{0}{场景 / 22}
\BookLine{139.00}{29.52}{}{0}{第一幕 / 23}
\BookLine{156.00}{29.52}{}{0}{第二幕 / 81}
\BookLine{173.00}{29.52}{}{0}{第三幕 / 157}
\BookLine{190.00}{29.52}{}{0}{醉酒的女神 / 227}
\newpage
% 原书页码：空白
\null
\newpage
% 原书页码：1
\null
\BookNumber{1}
\Vertical{15}{28}{}{14}{21}前\\言\EndVertical
\BookLine{105.00}{7.38}{}{0}{人们总是将生命看作一条沿着既定路线奔流的河，而在桑}
\BookLine{113.50}{0.00}{}{0}{顿·怀尔德的小说《第八日》的最后几页，我哥哥让其中一位}
\BookLine{122.00}{0.00}{}{0}{主角反驳了这种看法。相较之下，一个人的人生经历就像一张}
\BookLine{130.50}{0.00}{}{0}{织毯，它的长和宽都不确定，两面的纹路也不尽相同。彩色的}
\BookLine{139.00}{0.00}{}{0}{纱线勾勒出故事的轮廓，设计上也许极不对称，却清晰、真实}
\BookLine{147.50}{0.00}{}{0}{而完整。是记录，也是肖像。}
\BookLine{156.00}{7.38}{}{0}{桑顿的命运织毯起初还很小，七八岁时他生活在威斯康}
\BookLine{164.50}{0.00}{}{0}{星州的麦迪逊市，在阅读布尔凡奇的《寓言时代》时第一次听}
\BookLine{173.00}{0.00}{}{0}{到了阿尔刻提斯这个名字。在前基督教时期的希腊神话和歌}
\newpage
% 原书页码：2
\null
\BookNumber{2}
\Vertical{4}{169}{}{6.2}{8}阿\\尔\\刻\\提\\斯\\之\\歌\EndVertical
\BookLine{30.00}{0.00}{}{0}{谣中，国王佩利阿斯的女儿是公主中的公主。她的故事占据}
\BookLine{38.50}{0.00}{}{0}{了他的想象，也俘获了他的心灵。她的白色纱线嵌进了他的}
\BookLine{47.00}{0.00}{}{0}{命运织物中，阿尔刻提斯成为他潜意识中一份温柔的坚持，}
\BookLine{55.50}{0.00}{}{0}{一个有待发掘的灵感素材。她是一个戏剧雏形，直到多年后}
\BookLine{64.00}{0.00}{}{0}{才得以创作出来。}
\BookLine{72.50}{7.38}{}{0}{当怀尔德一家搬到加利福尼亚州伯克利时，阿尔刻提斯}
\BookLine{81.00}{0.00}{}{0}{之线打了个结。桉树林山坡上建有一座宏伟的希腊剧院，是}
\BookLine{89.50}{0.00}{}{0}{学校和小镇生活中生机勃勃的一部分。古典学系每年都会制作}
\BookLine{98.00}{0.00}{}{0}{几部索福克勒斯、埃斯库罗斯和欧里庇得斯的戏剧作品。我们}
\BookLine{106.50}{0.00}{}{0}{的母亲参与了服装组的志愿工作，用模板印出希腊钥匙的形}
\BookLine{115.00}{0.00}{}{0}{状，或是在漂亮的长袍上印上月桂树叶的图案。她给桑顿做了}
\BookLine{123.50}{0.00}{}{0}{一件有贝壳镶边的蓝色小袍子——给哥哥阿莫斯做了一件绿色}
\BookLine{132.00}{0.00}{}{0}{的——然后让他们去扮演雅典民众。桑顿就这样了解到了“整}
\BookLine{140.50}{0.00}{}{0}{个”剧场和古希腊的黄金时代，在这之前，他的全部剧场经验}
\BookLine{149.00}{0.00}{}{0}{不过是坐在密尔沃基剧院的后排看了一场《皆大欢喜》。}
\BookLine{157.50}{7.38}{}{0}{桑顿十岁了，黑头发，蓝眼睛，充满好奇心，精力旺盛。}
\BookLine{166.00}{0.00}{}{0}{早在十岁之前他就掌握了作家们的工具，驯服了二十六个字母}
\BookLine{174.50}{0.00}{}{0}{为自己服务，将它们组成词汇，在适当的时候用来表达自己。}
\newpage
% 原书页码：3
\null
\BookNumber{3}
\Vertical{129}{11}{}{7}{10}前\\言\EndVertical
\BookLine{30.00}{0.00}{}{0}{他早睡早起写作，随着词汇量的增长，他开始创造对话。他给}
\BookLine{38.50}{0.00}{}{0}{我们和邻居的孩子披上借来的纱，哄着我们一起朗诵他的豪言}
\BookLine{47.00}{0.00}{}{0}{壮语。}
\BookLine{55.50}{7.38}{}{0}{四年之后，他成为中国烟台一所教会学校的寄宿生。在一}
\BookLine{64.00}{0.00}{}{0}{所英国人开办的专为牛津、剑桥入学考试做准备的学校里，他}
\BookLine{72.50}{0.00}{}{0}{开始了对古典文学的学习。}
\BookLine{81.00}{7.38}{}{0}{他大学时代的前两年是在俄亥俄州的欧伯林学院度过的。}
\BookLine{89.50}{0.00}{}{0}{在那里他成了查尔斯·H. A. 瓦格教授的得意门生和朋友，而瓦}
\BookLine{98.00}{0.00}{}{0}{格教授本身也是一名出色的古典文学家。1920 年从耶鲁大学毕}
\BookLine{106.50}{0.00}{}{0}{业后，他在美国设立于罗马的古典学研究院待了八个月，阿尔}
\BookLine{115.00}{0.00}{}{0}{刻提斯之线在这些年里有了更清晰的纹路。尽管不是一名研究}
\BookLine{123.50}{0.00}{}{0}{员，他还是会去听讲座，和考古学家们一起去挖掘现场，仔细}
\BookLine{132.00}{0.00}{}{0}{研读莎草纸上的文字。在意大利的博物馆和美术馆里，他捕捉}
\BookLine{140.50}{0.00}{}{0}{到了阿尔刻提斯的影子。}
\BookLine{149.00}{7.38}{}{0}{在二十世纪二三十年代里，这张命运织毯延展得更大，也}
\BookLine{157.50}{0.00}{}{0}{变得更加密实。他开始从事教学工作——先是在劳伦斯维尔学}
\BookLine{166.00}{0.00}{}{0}{校教授法语，然后在芝加哥大学教授比较文学，其中也会讲到}
\BookLine{174.50}{0.00}{}{0}{古希腊的剧作家们——还有巡回演讲以及在好莱坞进行电影剧}
\newpage
% 原书页码：4
\null
\BookNumber{4}
\Vertical{4}{169}{}{6.2}{8}阿\\尔\\刻\\提\\斯\\之\\歌\EndVertical
\BookLine{30.00}{0.00}{}{0}{本创作。他出版了两部短剧集和四部小说，其中一部《圣路易}
\BookLine{38.50}{0.00}{}{0}{斯雷大桥》为他赢得了普利策奖，在三十岁的这一年使他名扬}
\BookLine{47.00}{0.00}{}{0}{四海。随之而来的是各种特权、荣誉和令人眼花缭乱的机遇，}
\BookLine{55.50}{0.00}{}{0}{同时他也失去了很多隐私，身体、思想和精神都受到了危害。}
\BookLine{64.00}{7.38}{}{0}{这几年里，桑顿的第一部长剧《我们的小镇》在制作演出}
\BookLine{72.50}{0.00}{}{0}{后便取得了成功。第二部剧《扬克斯商人》则以失败告终。安}
\BookLine{81.00}{0.00}{}{0}{德烈·奥比的剧作《卢克丽丝》由凯瑟琳·康奈尔委托给我哥}
\BookLine{89.50}{0.00}{}{0}{哥翻译，这部剧开演后又停演。与此同时，他帮鲁斯·戈登改}
\BookLine{98.00}{0.00}{}{0}{编的易卜生的《玩偶之家》则取得了巨大的成功。}
\BookLine{106.50}{7.38}{}{0}{1939 年夏天，桑顿终于完成了所有的邀约工作，带着完全}
\BookLine{115.00}{0.00}{}{0}{自由的状态和满腹个人创作计划回到家中。他已经写了两部作}
\BookLine{123.50}{0.00}{}{0}{品的开头，两部都是戏剧。一部还没有命名，另一部就是《阿}
\BookLine{132.00}{0.00}{}{0}{尔刻提斯之歌》。九月到了，希特勒的军队横扫了荷兰和比利}
\BookLine{140.50}{0.00}{}{1}{时，桑顿别无选择。《阿尔刻提斯之歌》并不适合一个战乱的世}
\BookLine{149.00}{0.00}{}{0}{界。此时尚未被命名的《九死一生》才更合适。}
\BookLine{157.50}{7.38}{}{0}{珍珠港事件后，桑顿被任命为空军上尉。很巧的是，他}
\BookLine{166.00}{0.00}{}{1}{出国服役前的最后一天，正是《九死一生》在纽约首演的日子。}
\BookLine{174.50}{0.00}{}{0}{我和母亲同他告别的那个早上，他拍了拍自己鼓鼓的行囊说}
\newpage
% 原书页码：5
\null
\BookNumber{5}
\Vertical{129}{11}{}{7}{10}前\\言\EndVertical
\BookLine{30.00}{0.00}{}{0}{道：“我打包了《阿尔刻提斯之歌》的手稿在里面。我一直在琢}
\BookLine{38.50}{0.00}{}{0}{磨把它写成一个不错的独幕剧，很可能只需要第二幕的那一大}
\BookLine{47.00}{0.00}{}{0}{场戏。我想留下一部完成的作品……说不定能派上用场，妈}
\BookLine{55.50}{0.00}{}{0}{妈。《九死一生》能不能成功也说不定呢……”}
\BookLine{64.00}{7.38}{}{0}{两年后，他弄丢了自己的手稿。彼时他已经是一位空军中}
\BookLine{72.50}{0.00}{}{0}{校，正在波士顿附近的军营里等待着退伍。指挥官批准他离开}
\BookLine{81.00}{0.00}{}{0}{三天，桑顿从后湾车站直奔科普利广场的波士顿图书馆，再次}
\BookLine{89.50}{0.00}{}{0}{沉浸到了古希腊的黄金时代之中。几天后，他终于在三年后又}
\BookLine{98.00}{0.00}{}{0}{一次坐在了康涅狄格汉姆登自己的书桌前。}
\BookLine{106.50}{7.38}{}{0}{此刻我面前正放着那时他留下的一张纸，一面用大写字母}
\BookLine{115.00}{0.00}{}{0}{打印着：}
\BookLine{132.00}{14.76}{\kaifont}{0}{《阿尔刻提斯之歌》}
\BookLine{140.50}{14.76}{\kaifont}{0}{一出充满问题的戏剧}
\BookLine{157.50}{7.38}{}{0}{反面是他写的一些小字：}
\BookLine{174.50}{14.76}{\kaifont}{0}{（第一幕戏的草稿，包括有忒瑞西阿斯的场景，都在}
\newpage
% 原书页码：6
\null
\BookNumber{6}
\Vertical{4}{169}{}{6.2}{8}阿\\尔\\刻\\提\\斯\\之\\歌\EndVertical
\BookLine{30.00}{7.38}{\kaifont}{0}{战争前完成，已经遗失）}
\BookLine{38.50}{14.76}{\kaifont}{0}{第一幕初稿，1945 年 5 月和 7 月}
\BookLine{47.00}{14.76}{\kaifont}{0}{第二稿，1945 年 7 月 8 日开始创作}
\BookLine{55.50}{14.76}{\kaifont}{0}{第三稿，12 月开始（在第二幕与第三幕的一半完}
\BookLine{64.00}{7.38}{\kaifont}{0}{成后）}
\BookLine{81.00}{7.38}{}{0}{没写上的日期应该是 1945 年。}
\BookLine{89.50}{7.38}{}{0}{遗失手稿给哥哥带来的影响远不是这几行笔记可以传达}
\BookLine{98.00}{0.00}{}{0}{的。他急切地想要回到和平时期的状态中，这激发了他的热}
\BookLine{106.50}{0.00}{}{0}{情，让他如记录中所写的样子，连续创作了七个月的时间。然}
\BookLine{115.00}{0.00}{}{0}{而，就如同其他上百万的士兵们一样，他回到家时不只是迷茫}
\BookLine{123.50}{0.00}{}{0}{与疲惫，整个人也完全改变了。他不再能够以战争之前那样的}
\BookLine{132.00}{0.00}{}{0}{视角来看待阿尔刻提斯了，他的体力和精力也不足以支撑自己}
\BookLine{140.50}{0.00}{}{0}{保持一个亢奋的状态来完成工作。他发现自己所做的更多是破}
\BookLine{149.00}{0.00}{}{0}{坏而不是创造。好在他意识到了这一点，在造成破坏前就停止}
\BookLine{157.50}{0.00}{}{0}{了自己的工作。}
\BookLine{166.00}{7.38}{}{0}{古代文学故事里满是战场和战士，暴君和英雄。哥哥的思}
\BookLine{174.50}{0.00}{}{0}{绪还围绕着他自己在希特勒和墨索里尼带来的战争中的经历：}
\newpage
% 原书页码：7
\null
\BookNumber{7}
\Vertical{129}{11}{}{7}{10}前\\言\EndVertical
\BookLine{30.00}{0.00}{}{0}{他在北非待了一年，就在那些希腊与罗马帝国留下的废墟与阴}
\BookLine{38.50}{0.00}{}{0}{影里；又在意大利南部待了一年，那里同样充满戏剧性的传}
\BookLine{47.00}{0.00}{}{0}{说，例如恺撒的胜利。桑顿突然有了新作品的灵感，他的小说}
\BookLine{55.50}{0.00}{}{1}{《三月十五》在 1948 年出版。这部小说开启了他回归文明世界}
\BookLine{64.00}{0.00}{}{0}{的道路，他又开始教书，演讲，上剧场，重新思考阿尔刻提斯}
\BookLine{72.50}{0.00}{}{0}{的故事。}
\BookLine{81.00}{7.38}{}{0}{在欧里庇得斯之后，有许多作家想要以自己的方式改}
\BookLine{89.50}{0.00}{}{0}{编阿尔刻提斯的故事。1773 年，克里斯托夫·马丁·维兰德}
\BookLine{98.00}{0.00}{}{0}{以此为题创作了一部歌剧剧本。在十九世纪，奥地利诗人胡}
\BookLine{106.50}{0.00}{}{0}{戈·冯·霍夫曼斯塔尔和英国作家罗伯特·勃朗宁也都选取过}
\BookLine{115.00}{0.00}{}{1}{这个题材。T. S. 艾略特在他成功的剧作《鸡尾酒会》中也有所借}
\BookLine{123.50}{0.00}{}{0}{用。也难怪，这个故事的主旨是夫妻之间的爱，以及妻子愿意}
\BookLine{132.00}{0.00}{}{0}{为了丈夫而死。这对任何想要发掘和赞美人类能动性的作者来}
\BookLine{140.50}{0.00}{}{0}{说，都是充满吸引力且发人深省的挑战。}
\BookLine{149.00}{7.38}{}{0}{1954 年，《媒人》在爱丁堡艺术节上演并获得成功——这}
\BookLine{157.50}{0.00}{}{0}{部有趣的作品由纽约的失败之作《扬克斯商人》而来，后来发}
\BookLine{166.00}{0.00}{}{0}{展成为大热的音乐剧《你好，多莉！》——制作人委托我哥哥}
\BookLine{174.50}{0.00}{}{0}{为下一季创作新作品，这给了他完成《阿尔刻提斯之歌》的动}
\newpage
% 原书页码：8
\null
\BookNumber{8}
\Vertical{4}{169}{}{6.2}{8}阿\\尔\\刻\\提\\斯\\之\\歌\EndVertical
\BookLine{30.00}{0.00}{}{0}{力。他在普罗旺斯的艾克斯闭关创作了很长一段时间。圣诞}
\BookLine{38.50}{0.00}{}{0}{节时，他聚集了包括我在内的一帮朋友，一起在瑞士的维拉}
\BookLine{47.00}{0.00}{}{0}{斯待了两周。在此期间，我依据他的手稿打出了新文本的前}
\BookLine{55.50}{0.00}{}{0}{两幕，他立刻提交给了艺术节组委会。泰隆·格思里接受这}
\BookLine{64.00}{0.00}{}{0}{出戏参与到 1955 年的节日中，春天离开艾克斯之前，他提交}
\BookLine{72.50}{0.00}{}{0}{了第三幕。}
\BookLine{81.00}{7.38}{}{0}{古希腊剧场的演出由内容相关联的悲剧三部曲组成，通常}
\BookLine{89.50}{0.00}{}{0}{还有一出羊人剧\textsuperscript{1}——用前面剧中的人物和情节进行或幽默或尖}
\BookLine{98.00}{0.00}{}{0}{锐的注解。欧里庇得斯是三大对手剧作家中最年轻的一位，他}
\BookLine{106.50}{0.00}{}{0}{有时会打破传统，写一部可以独立上演的剧本，或者该剧可以}
\BookLine{115.00}{0.00}{}{0}{与其他剧目组合成为三部曲。他的《阿尔刻提斯》就是这个类}
\BookLine{123.50}{0.00}{}{0}{型的剧目之一。事实上，由于它的“大团圆结局”——女主角}
\BookLine{132.00}{0.00}{}{1}{被赫拉克勒斯从冥界带回人间，《阿尔刻提斯》偶尔甚至会被用}
\BookLine{140.50}{0.00}{}{0}{作羊人剧演出。}
\BookLine{149.00}{7.38}{}{0}{我哥哥对于他这版改编的第一个想法，是将欧里庇得斯的}
\BookLine{157.50}{0.00}{}{0}{文本进行全新翻译作为中间一幕，然后创作第一幕和第三幕把}
\begin{textblock*}{45mm}(16mm,171mm)\rule{43mm}{.3pt}\end{textblock*}
\begin{textblock*}{108mm}(16mm,177mm)\kaifont\fontsize{7}{10}\selectfont 1　羊人剧，是一种轻松的笑剧，不是喜剧。一般在三部悲剧上演后进行演出，\end{textblock*}
\begin{textblock*}{108mm}(16mm,182mm)\kaifont\fontsize{7}{10}\selectfont 　 作为一种调剂。——译者注\end{textblock*}
\newpage
% 原书页码：9
\null
\BookNumber{9}
\Vertical{129}{11}{}{7}{10}前\\言\EndVertical
\BookLine{30.00}{0.00}{}{0}{第二幕框起来。很快他就放弃了这个计划，因为这不符合他的}
\BookLine{38.50}{0.00}{}{0}{目的。}
\BookLine{47.00}{7.38}{}{0}{在希腊人讲述的传说和几个世纪后布尔凡奇的记录中，阿}
\BookLine{55.50}{0.00}{}{0}{德墨托斯国王知道，如果有人愿意为他而死，他就可以活下}
\BookLine{64.00}{0.00}{}{0}{去。他向他的仆人、昔日的军中同僚和朋友们求救，但无济于}
\BookLine{72.50}{0.00}{}{0}{事。而后，他不遗余力地恳求他的父母，指出既然他们的寿命}
\BookLine{81.00}{0.00}{}{0}{所剩无几，请求他们代死只是小事一桩。但无论是老父亲还是}
\BookLine{89.50}{0.00}{}{0}{老母亲，都不愿意为儿子放弃生命。国王没有请求他的王后做}
\BookLine{98.00}{0.00}{}{0}{出这个牺牲，当阿尔刻提斯带着爱献出她的生命时，他的抗议}
\BookLine{106.50}{0.00}{}{0}{显得不温不火，毫无意义。}
\BookLine{115.00}{7.38}{}{0}{桑顿让阿德墨托斯不曾得知自己可以通过他人的死亡得到}
\BookLine{123.50}{0.00}{}{0}{拯救，在第二幕造就了一个戏剧性和情感的高潮，并深化了人}
\BookLine{132.00}{0.00}{}{0}{物形象和夫妻关系，为这个传说原本简单的故事线赋予了来自}
\BookLine{140.50}{0.00}{}{0}{二十世纪的丰富解读。关于他的这个版本，哥哥说过：}
\BookLine{157.50}{14.76}{\kaifont}{0}{从某个层面上来说，我的剧本描述了一个女人——}
\BookLine{166.00}{7.38}{\kaifont}{0}{许多女人——从一个犹疑的新娘成长为经受着苦考验的}
\BookLine{174.50}{7.38}{\kaifont}{0}{妻子，最后成为不堪重负的老人。在另一个层面上，它}
\newpage
% 原书页码：10
\null
\BookNumber{10}
\Vertical{4}{169}{}{6.2}{8}阿\\尔\\刻\\提\\斯\\之\\歌\EndVertical
\BookLine{30.00}{7.38}{\kaifont}{0}{是一个关于神和人，关于死亡和冥界，关于伟大的爱和}
\BookLine{38.50}{7.38}{\kaifont}{0}{伟大的试炼后的复活，关于篡夺和复仇的疯狂而浪漫的}
\BookLine{47.00}{7.38}{\kaifont}{0}{故事。在又一个层面上，它是一部喜剧……关于天界和}
\BookLine{55.50}{7.38}{\kaifont}{0}{人间沟通上的极端困难，关于人类和神明的不可比性所}
\BookLine{64.00}{7.38}{\kaifont}{0}{造成的误解。\textsuperscript{1}}
\BookLine{81.00}{7.38}{}{0}{在桑顿的其他作品中，1930 年出版的第三部小说《安德罗}
\BookLine{89.50}{0.00}{}{1}{斯的女人》在设定、情感和思想上与《阿尔刻提斯之歌》最为接}
\BookLine{98.00}{0.00}{}{0}{近。故事发生在希腊的一个小岛上，彼时阿尔刻提斯已经成为}
\BookLine{106.50}{0.00}{}{0}{传说。这部作品是一首怀旧的、宁静却也充满激情的散文诗，}
\BookLine{115.00}{0.00}{}{0}{它像一座桥，一端陷在残酷的异教世界的泥沼中，另一端延伸}
\BookLine{123.50}{0.00}{}{0}{向正缓缓到来的解放的黎明。女主角克里西丝漂亮、聪明，受}
\BookLine{132.00}{0.00}{}{0}{过一定的教育，但因为自己交际花的身份被边缘化。然而，她}
\BookLine{140.50}{0.00}{}{0}{是阿尔刻提斯女王精神上的姐妹。}
\BookLine{149.00}{7.38}{}{0}{通过这两个女性角色，桑顿向所有女性表达了敬意，也让}
\BookLine{157.50}{0.00}{}{0}{角色承担起传递自己主旨的重任。克里西丝在她的生和死中都}
\begin{textblock*}{45mm}(16mm,171mm)\rule{43mm}{.3pt}\end{textblock*}
\begin{textblock*}{108mm}(16mm,177mm)\kaifont\fontsize{7}{10}\selectfont 1　引自保罗·M. 库贝塔的《现代戏剧分析》（纽约：霍尔特，莱因哈特和温斯顿公司，\end{textblock*}
\begin{textblock*}{108mm}(16mm,182mm)\kaifont\fontsize{7}{10}\selectfont 　 1955），第 595—596 页。\end{textblock*}
\newpage
% 原书页码：11
\null
\BookNumber{11}
\Vertical{129}{11}{}{7}{10}前\\言\EndVertical
\BookLine{30.00}{0.00}{}{0}{赞颂了全部生命——“光明抑或黑暗”——并教导其他人也要}
\BookLine{38.50}{0.00}{}{0}{这样做。她在不自知的情形下成了一位先知，为人类指明了更}
\BookLine{47.00}{0.00}{}{0}{高的希望所在。}
\BookLine{55.50}{7.38}{}{0}{当然，阿尔刻提斯这个名字，早在桑顿听说她之前，就}
\BookLine{64.00}{0.00}{}{0}{已经承载了许多意义，被月桂花环围绕。在他的剧中，阿尔刻}
\BookLine{72.50}{0.00}{}{0}{提斯从冥界被带回，重新生活在阿波罗的光芒之中，她的结局}
\BookLine{81.00}{0.00}{}{0}{也与人类不同。神把她从死神那里偷了回来，坚持要她永远活}
\BookLine{89.50}{0.00}{}{0}{着。阿波罗将她从阿德墨托斯身边偷走，也夺去了她——以及}
\BookLine{98.00}{0.00}{}{0}{国王为她在皇家陵墓中立碑的机会。阿波罗把她带到自己的长}
\BookLine{106.50}{0.00}{}{0}{青林中，她仍在那里徘徊，成为一个象征，一个神话，一个讲}
\BookLine{115.00}{0.00}{}{0}{给我们所有人的真理：在基督教人文主义精神看来，一个逃离}
\BookLine{123.50}{0.00}{}{0}{了黑暗的人，她自己就是地平线上的一束光。}
\BookLine{132.00}{7.38}{}{0}{一部以古希腊戏剧形式写下的戏剧，理论上应该以传统}
\BookLine{140.50}{0.00}{}{0}{的古典形式在一个大理石建造的露天剧场中演出。桑顿创作}
\BookLine{149.00}{0.00}{}{0}{这版《阿尔刻提斯之歌》，也是为了它能够适应传统的镜框式}
\BookLine{157.50}{0.00}{}{0}{舞台。在爱丁堡，这出戏的演出环境有所不同，是在一个叫}
\BookLine{166.00}{0.00}{}{0}{作苏格兰教会大会堂的巨大谷仓里进行表演的。剧院呈箱形，}
\BookLine{174.50}{0.00}{}{0}{一排排并不舒适的座位构成了一个三面舞台——若是要为}
\newpage
% 原书页码：12
\null
\BookNumber{12}
\Vertical{4}{169}{}{6.2}{8}阿\\尔\\刻\\提\\斯\\之\\歌\EndVertical
\BookLine{30.00}{0.00}{}{0}{教会政治和教会条例而在会议上互相争斗，这倒是个完美的}
\BookLine{38.50}{0.00}{}{0}{环境。}
\BookLine{47.00}{7.38}{}{1}{在这个不大合拍的演出环境中，《阿尔刻提斯之歌》这个简}
\BookLine{55.50}{0.00}{}{0}{单的故事不得不被放大到一个夸张的比例。泰隆·格思里导演}
\BookLine{64.00}{0.00}{}{0}{的处理激动人心却过度戏剧化了，他增加了舞台行动，用群戏}
\BookLine{72.50}{0.00}{}{0}{场面填满舞台和过道。且不说是好是坏，这出戏的本质完全丧}
\BookLine{81.00}{0.00}{}{0}{失了。}
\BookLine{89.50}{7.38}{}{0}{甚至连剧名都改成了《阳光下的一生》，委员会认为这个}
\BookLine{98.00}{0.00}{}{0}{名字“更具票房吸引力”。桑顿从苏格兰给身在麻省剑桥的哥}
\BookLine{106.50}{0.00}{}{0}{哥阿莫斯写的信中提到了这一点：}
\BookLine{123.50}{14.76}{\kaifont}{0}{这就是《阿尔刻提斯之歌》，但老板们不想用这个名}
\BookLine{132.00}{7.38}{\kaifont}{0}{字。对于一出充满了可怖情节的戏来说，《阳光下的一}
\BookLine{140.50}{7.38}{\kaifont}{0}{生》这个名字过于轻松愉快了，但他们说阿尔刻提斯的}
\BookLine{149.00}{7.38}{\kaifont}{0}{生命源于阿波罗，而阿波罗就是太阳，所以就有了这个}
\BookLine{157.50}{7.38}{\kaifont}{0}{名字……}
\BookLine{174.50}{7.38}{}{0}{桑顿还是得到了一些补偿。他不是一个会绝望的人，甚}
\newpage
% 原书页码：13
\null
\BookNumber{13}
\Vertical{129}{11}{}{7}{10}前\\言\EndVertical
\BookLine{30.00}{0.00}{}{0}{至也不会焦虑太久。格思里的才华创造了一些精彩段落，以艾}
\BookLine{38.50}{0.00}{}{0}{琳·沃思为首的是一个强大的演员阵容。尽管批评家们还是老}
\BookLine{47.00}{0.00}{}{0}{样子，先是充满兴趣与理解，随即便是因全知而不屑。而大部}
\BookLine{55.50}{0.00}{}{0}{分的观众尽管坐在很不舒服的位置上，但都给出了积极的反响}
\BookLine{64.00}{0.00}{}{0}{和明确的肯定。}
\BookLine{72.50}{7.38}{}{0}{最重要的是，作者看到自己的作品立体而生动地呈现了}
\BookLine{81.00}{0.00}{}{0}{出来，他知道自己做到了哪些，又有哪些不足。当泰隆·格}
\BookLine{89.50}{0.00}{}{0}{思里非常有雅量地向他吐露“我没有把你的戏演好”时，桑}
\BookLine{98.00}{0.00}{}{0}{顿也得到了安慰，他觉得是被苏格兰教会大会堂这个场地给}
\BookLine{106.50}{0.00}{}{0}{坑了。}
\BookLine{115.00}{7.38}{}{0}{国内外有许多戏剧制作人对这部作品感兴趣，想获得授}
\BookLine{123.50}{0.00}{}{0}{权在纽约制作。哥哥对于他想要在文本上进行的调整也有自}
\BookLine{132.00}{0.00}{}{0}{己的想法。他与一些导演进行了交流，听取了各种建议，开始}
\BookLine{140.50}{0.00}{}{0}{着手修改。这对他来说一直不是件容易的事，他擅长的是活在}
\BookLine{149.00}{0.00}{}{0}{当下，着眼于未来。有太多人围着他给他建议，于是他告诉他}
\BookLine{157.50}{0.00}{}{0}{们：“我需要更长远的考量。”他收回了剧本，拒绝了那些想要}
\BookLine{166.00}{0.00}{}{0}{基于此版本制作的请求。}
\BookLine{174.50}{7.38}{}{0}{当他从这种完全无法适应的压力中解脱出来后，他开始有}
\newpage
% 原书页码：14
\null
\BookNumber{14}
\Vertical{4}{169}{}{6.2}{8}阿\\尔\\刻\\提\\斯\\之\\歌\EndVertical
\BookLine{30.00}{0.00}{}{0}{了新的想法，想要按照希腊戏剧的传统来创作《阿尔刻提斯之}
\BookLine{38.50}{0.00}{}{1}{歌》和羊人剧。很快，他就沉浸在《醉酒的女神》的创作中了。}
\BookLine{47.00}{0.00}{}{0}{这是作家或各种形式的创作者偶尔会经历的状况，他一下就}
\BookLine{55.50}{0.00}{}{0}{“想到”了这部作品，就像一条直线，从头到尾一气呵成。作}
\BookLine{64.00}{0.00}{}{1}{品 1957 年在《大西洋月刊》上发表，和原稿相比几乎没有任何}
\BookLine{72.50}{0.00}{}{0}{改动。}
\BookLine{81.00}{7.38}{}{0}{因为希望《醉酒的女神》能在舞台上呈现，我哥哥签了合}
\BookLine{89.50}{0.00}{}{0}{约，让这两部戏能够在瑞士苏黎世的剧院上演。第二次世界}
\BookLine{98.00}{0.00}{}{1}{大战期间，《我们的小镇》和《九死一生》都是在这个剧院进行}
\BookLine{106.50}{0.00}{}{0}{的欧洲首演。德国和奥地利最知名的导演和演员们都在瑞士避}
\BookLine{115.00}{0.00}{}{0}{难，桑顿的叙事作品和戏剧作品，都为他们所熟知并给予了高}
\BookLine{123.50}{0.00}{}{0}{度评价。}
\BookLine{132.00}{7.38}{}{0}{桑顿与译者赫伯特·E. 赫利奇卡合作，赫利奇卡是维也纳}
\BookLine{140.50}{0.00}{}{0}{人，曾将他早期的许多作品翻译成德文。相较爱丁堡的演出版}
\BookLine{149.00}{0.00}{}{0}{本，桑顿做的修改并不大，他希望看到另一位导演和另一个团}
\BookLine{157.50}{0.00}{}{0}{队来呈现同一个剧本。这个版本的制作在各种意义上都令人激}
\BookLine{166.00}{0.00}{}{0}{动和震撼，观众在剧院里高度专注，在各幕之间进行讨论，演}
\BookLine{174.50}{0.00}{}{0}{出结束时用力鼓掌，离开剧院时活力满满。}
\newpage
% 原书页码：15
\null
\BookNumber{15}
\Vertical{129}{11}{}{7}{10}前\\言\EndVertical
\BookLine{30.00}{7.38}{}{0}{苏黎世首演之后，该剧在全德国大大小小城市的国有剧}
\BookLine{38.50}{0.00}{}{0}{院都进行了完整版的演出，维也纳城堡剧院也上演了该剧。在}
\BookLine{47.00}{0.00}{}{0}{此期间，桑顿做了进一步的修改，为了回应在德国的巨大反}
\BookLine{55.50}{0.00}{}{1}{响，《阿尔刻提斯之歌》在 1960 年由费舍尔出版社出版了德文}
\BookLine{64.00}{0.00}{}{0}{版剧本。}
\BookLine{72.50}{7.38}{}{0}{该剧在报纸和杂志的文学栏目，以及学术期刊和书籍中都}
\BookLine{81.00}{0.00}{}{0}{引起了讨论。凯特·汉伯格在她关于现代文学中的古典人物的}
\BookLine{89.50}{0.00}{}{0}{集子《从索福克勒斯到萨特》\textsuperscript{1}中写道：“怀尔德的作品是现代世}
\BookLine{98.00}{0.00}{}{0}{界文学中对阿尔刻提斯这一题材最重要的表达。”}
\BookLine{106.50}{7.38}{}{0}{尽管桑顿没有继续把《阿尔刻提斯之歌》写成一个英文的}
\BookLine{115.00}{0.00}{}{0}{戏剧剧本，他也没有放任不管。他在为自己非常欣赏的作曲家}
\BookLine{123.50}{0.00}{}{0}{露易丝·塔尔玛寻找歌剧剧本题材时，似乎不可避免地完成了}
\BookLine{132.00}{0.00}{}{0}{这个循环。}
\BookLine{140.50}{7.38}{}{0}{这部歌剧也叫《阿尔刻提斯之歌》，历时六年完成制作，}
\BookLine{149.00}{0.00}{}{1}{1962 年 3 月 2 日在德国法兰克福重建后的极为华丽的歌剧院进}
\BookLine{157.50}{0.00}{}{0}{行了首演。这是一部规模宏大的歌剧，需要动用很大的舞台、}
\begin{textblock*}{45mm}(16mm,171mm)\rule{43mm}{.3pt}\end{textblock*}
\begin{textblock*}{108mm}(16mm,177mm)\kaifont\fontsize{7}{10}\selectfont 1　纽约：昂格尔，1969。编译自 Von Sophokles zu Sartre, Griechische Dramenfiguren\end{textblock*}
\begin{textblock*}{108mm}(16mm,182mm)\kaifont\fontsize{7}{10}\selectfont 　 Antik und Modern, Stuttgart, 1962.\end{textblock*}
\newpage
% 原书页码：16
\null
\BookNumber{16}
\Vertical{4}{169}{}{6.2}{8}阿\\尔\\刻\\提\\斯\\之\\歌\EndVertical
\BookLine{30.00}{0.00}{}{0}{大规模管弦乐队和一大批演员。在补贴的支持下，制作得以满}
\BookLine{38.50}{0.00}{}{0}{足这些要求，英格·博克将阿尔刻提斯演绎得极为出色。大幕}
\BookLine{47.00}{0.00}{}{0}{落下时，剧院里响起雷鸣般的掌声和此起彼伏的叫好声。大家}
\BookLine{55.50}{0.00}{}{0}{进行了三十多次谢幕，塔尔玛小姐、桑顿，每个人都被一次又}
\BookLine{64.00}{0.00}{}{0}{一次地叫上去鞠躬致意。演出结束后的晚宴上，大家都十分雀}
\BookLine{72.50}{0.00}{}{0}{跃。不过，在接下来的几天里和后来的周刊中，评论家们的反}
\BookLine{81.00}{0.00}{}{0}{馈褒贬不一，其中一个隐晦但明确的抱怨是女人不应该写歌}
\BookLine{89.50}{0.00}{}{0}{剧，还有就是音乐过于现代了。}
\BookLine{98.00}{7.38}{}{0}{对桑顿而言，他对配乐完全满意，他认为作品的不足之处}
\BookLine{106.50}{0.00}{}{0}{都源于自己作为词作者的经验不足。这部歌剧后来常以音乐会}
\BookLine{115.00}{0.00}{}{0}{的形式上演，其中的咏叹调也在各种独唱会中被演唱。桑顿从}
\BookLine{123.50}{0.00}{}{0}{未怀疑过，这是一部领先于时代的作品，总有一天会得到充分}
\BookLine{132.00}{0.00}{}{0}{的认可。}
\BookLine{140.50}{7.38}{}{0}{虽然我哥哥很早就开始写作，但他的作品列表并不长：七}
\BookLine{149.00}{0.00}{}{0}{部小说，一些论说性散文和关于自己与他人作品的序言，两部}
\BookLine{157.50}{0.00}{}{1}{短剧集和四部长剧，其中《扬克斯商人》和《媒人》算作一部。}
\BookLine{166.00}{0.00}{}{0}{在作品中，他不怕重复自己，也不怕受到薪火相传的传统大师}
\BookLine{174.50}{0.00}{}{0}{们的影响。}
\newpage
% 原书页码：17
\null
\BookNumber{17}
\Vertical{129}{11}{}{7}{10}前\\言\EndVertical
\BookLine{30.00}{7.38}{}{1}{关于他的戏剧作品，他在集子《戏剧三种》\textsuperscript{1}的序言中写道：}
\BookLine{47.00}{14.76}{\kaifont}{0}{我不是大家所期待的创新剧作家之一。我希望我是。}
\BookLine{55.50}{7.38}{\kaifont}{0}{我希望自己的创作是为他们铺路。我不是一个创新者，而}
\BookLine{64.00}{7.38}{\kaifont}{0}{是重新发现那些被遗忘的事物的人，希望我能清除那些妨}
\BookLine{72.50}{7.38}{\kaifont}{0}{碍着我们的东西。}
\BookLine{89.50}{7.38}{}{0}{但他对形式进行了实验，敢于对现实时间、历史时间和文}
\BookLine{98.00}{0.00}{}{0}{学时间进行创造，从历史中摘取事件，从故事和绘画中摘取人}
\BookLine{106.50}{0.00}{}{0}{物，从神圣传说中摘取角色。他让死亡、天使、异教神灵和动}
\BookLine{115.00}{0.00}{}{0}{物发声。}
\BookLine{123.50}{7.38}{}{0}{在他的三分钟剧本集《搅乱水池的天使及其他作品》\textsuperscript{2}的前}
\BookLine{132.00}{0.00}{}{0}{言中，我哥哥写到一个年轻作家的问题{\kaifont （他了解这些问题，因}}
\BookLine{140.50}{0.00}{}{0}{{\kaifont 为他在十几岁的时候就写了集子中的一些剧本）}。他描述了写}
\BookLine{149.00}{0.00}{}{0}{作时欣喜和沮丧的时刻，以及各种各样的可能性：}
\begin{textblock*}{45mm}(16mm,171mm)\rule{43mm}{.3pt}\end{textblock*}
\begin{textblock*}{108mm}(16mm,177mm)\kaifont\fontsize{7}{10}\selectfont 1　纽约：哈珀兄弟出版公司，1957。\end{textblock*}
\begin{textblock*}{108mm}(16mm,182mm)\kaifont\fontsize{7}{10}\selectfont 2　纽约：科沃德－麦肯出版公司，1928。\end{textblock*}
\newpage
% 原书页码：18
\null
\BookNumber{18}
\Vertical{4}{169}{}{6.2}{8}阿\\尔\\刻\\提\\斯\\之\\歌\EndVertical
\BookLine{30.00}{14.76}{\kaifont}{0}{如果一个人不允许自己自命不凡，那么写作的实践将}
\BookLine{38.50}{7.38}{\kaifont}{0}{会多么不同。有些人别无选择：他们宁愿写一部失败的清}
\BookLine{47.00}{7.38}{\kaifont}{0}{唱剧，也不愿写一首成功的民谣曲。}
\BookLine{64.00}{7.38}{}{0}{桑顿始终坚信，废纸篓是作家最好的朋友。即便如此，现}
\BookLine{72.50}{0.00}{}{0}{存的文档还是证明他写过一些“失败的清唱剧”，包括只有前}
\BookLine{81.00}{0.00}{}{0}{两幕，第三幕从没创作出来的《百货商场》。这部作品以世界}
\BookLine{89.50}{0.00}{}{0}{上最大、最诱人的商店为背景，在那里你可以找到你需要或想}
\BookLine{98.00}{0.00}{}{0}{要的一切。还有一部未完成的史诗剧情片剧本《熔炉》，通过}
\BookLine{106.50}{0.00}{}{0}{几代人的生活反映我们国家的发展。尽管他没有从其他项目中}
\BookLine{115.00}{0.00}{}{0}{抽出时间来修改《阿尔刻提斯之歌》的英文剧本，使其与德文}
\BookLine{123.50}{0.00}{}{0}{版内容保持一致，而且他肯定还会对剧本做出进一步调整，但}
\BookLine{132.00}{0.00}{}{0}{目前的剧本并不算是“失败的清唱剧”。这部作品成功地反映}
\BookLine{140.50}{0.00}{}{0}{了他对人类经验真实性最持久的信念。}
\BookLine{149.00}{7.38}{}{0}{桑顿·怀尔德的织毯已经完成。随着《阿尔刻提斯之歌》}
\BookLine{157.50}{0.00}{}{0}{与《醉酒的女神》的出版和故事的讲述，他的全部设计得以清}
\BookLine{166.00}{0.00}{}{0}{晰展现出来，以自己的文字能力，以自己的天赋，以人生的全}
\BookLine{174.50}{0.00}{}{0}{部岁月，最重要的，以自身的信念和分享的愿望。现在，阿}
\newpage
% 原书页码：19
\null
\BookNumber{19}
\Vertical{129}{11}{}{7}{10}前\\言\EndVertical
\BookLine{30.00}{0.00}{}{0}{尔刻提斯这团巨大的纱线，最后一段也已经解开，最后一针也}
\BookLine{38.50}{0.00}{}{0}{打好了结。在这个故事里，已经再没有谜团需要德尔斐来守}
\BookLine{47.00}{0.00}{}{0}{护了。}
\BookLine{64.00}{66.43}{}{0}{伊莎贝尔·怀尔德\textsuperscript{1}}
\BookLine{72.50}{81.19}{}{0}{玛莎葡萄园}
\BookLine{81.00}{81.19}{}{1}{1977 年 7 月}
\begin{textblock*}{45mm}(16mm,171mm)\rule{43mm}{.3pt}\end{textblock*}
\begin{textblock*}{108mm}(16mm,177mm)\kaifont\fontsize{7}{10}\selectfont 1　作者为桑顿·怀尔德的妹妹，她是一位小说家和传记作家，同时也是桑顿·怀\end{textblock*}
\begin{textblock*}{108mm}(16mm,182mm)\kaifont\fontsize{7}{10}\selectfont 　 尔德的文学经纪人、发言人和传记作者。——译者注\end{textblock*}
\newpage
% 原书页码：20
\null
\Vertical{15}{28}{}{14}{21}人\\物\EndVertical
\BookLine{82.00}{0.00}{}{0}{阿波罗——太阳神}
\BookLine{90.50}{0.00}{}{0}{死神——冥界之神}
\BookLine{99.00}{0.00}{}{0}{第一个看门人——守夜人，中年，阿德墨托斯国王的仆人}
\BookLine{107.50}{0.00}{}{0}{阿尔刻提斯——塞萨利王后}
\BookLine{116.00}{0.00}{}{0}{阿格莱亚——年迈的女仆}
\BookLine{124.50}{0.00}{}{0}{忒瑞西阿斯——盲人，年纪非常大，暴躁易怒，控制欲强，}
\BookLine{133.00}{25.83}{}{0}{马上就要老糊涂的状态}
\BookLine{141.50}{0.00}{}{0}{男孩——忒瑞西阿斯的仆人}
\BookLine{150.00}{0.00}{}{0}{牧人一——牧羊人}
\BookLine{158.50}{0.00}{}{0}{牧人二——牧羊人}
\BookLine{167.00}{0.00}{}{0}{牧人三——牧羊人}
\BookLine{175.50}{0.00}{}{0}{牧人四——牧羊人}
\BookLine{184.00}{0.00}{}{0}{罗德普——王宫中年轻的女仆}
\newpage
% 原书页码：21
\null
\BookNumber{21}
\BookLine{88.00}{0.00}{}{0}{赫拉克勒斯——宙斯之子}
\BookLine{96.50}{0.00}{}{0}{第二个看门人——比第一个看门人年轻，阿吉斯国王的仆人}
\BookLine{105.00}{0.00}{}{0}{伊皮伊尼兹——21 岁，阿尔刻提斯的儿子}
\BookLine{113.50}{0.00}{}{0}{切瑞安德——伊皮伊尼兹的朋友}
\BookLine{122.00}{0.00}{}{0}{阿吉斯——篡位的塞萨利国王}
\BookLine{130.50}{0.00}{}{0}{守卫一——阿吉斯王宫的守卫}
\BookLine{139.00}{0.00}{}{0}{守卫二——阿吉斯王宫的守卫}
\BookLine{147.50}{0.00}{}{0}{守卫三——阿吉斯王宫的守卫}
\BookLine{156.00}{0.00}{}{0}{守卫四——阿吉斯王宫的守卫}
\BookLine{164.50}{0.00}{}{0}{仆人们——阿德墨托斯国王和阿吉斯国王的仆人们}
\BookLine{173.00}{0.00}{}{0}{塞萨利的百姓}
\newpage
% 原书页码：22
\null
\Vertical{15}{28}{}{14}{21}场\\景\EndVertical
\BookLine{82.00}{0.00}{}{0}{塞萨利国王阿德墨托斯王宫的后院，}
\BookLine{90.50}{0.00}{}{0}{希腊黄金时代几个世纪之前。}
\newpage
% 原书页码：23
\null
\ActOrnament
\newpage
% 原书页码：24
\null
\newpage
% 原书页码：25
\null
\BookNumber{25}
\Vertical{129}{11}{}{7}{10}第\\一\\幕\EndVertical
\BookLine{30.00}{7.38}{\kaifont}{0}{除了第三幕结尾和《醉酒的女神》结束后，全剧不设大}
\BookLine{38.50}{0.00}{\kaifont}{0}{幕。《阿尔刻提斯之歌》全剧三幕戏的场景均在塞萨利国王阿}
\BookLine{47.00}{0.00}{\kaifont}{0}{德墨托斯宫殿的后院，时间是希腊黄金时代几个世纪之前。}
\BookLine{55.50}{0.00}{\kaifont}{0}{每一幕戏都开始于一天的黎明，结束于同一日的黄昏。}
\BookLine{64.00}{7.38}{\kaifont}{0}{宫殿是一座由粗粝石块砌成的低矮却宽敞的房屋，有着}
\BookLine{72.50}{0.00}{\kaifont}{0}{平坦的屋顶。巨大的树干仿佛组成了宫殿的门廊。宫殿大门}
\BookLine{81.00}{0.00}{\kaifont}{0}{是木制的，两扇门上各镶有一块镀金的牛骨。门前是平台，}
\BookLine{89.50}{0.00}{\kaifont}{0}{下几级台阶就能走到被松软土壤覆盖的庭院。}
\BookLine{98.00}{7.38}{\kaifont}{0}{宫殿的前区占据了舞台后方左侧四分之三的空间，舞台}
\BookLine{106.50}{0.00}{\kaifont}{0}{其他区域则被泥砖墙围了起来。右边墙上是一扇大木门，直}
\BookLine{115.00}{0.00}{\kaifont}{0}{通外面的大路和费莱城。左边的墙稍低一些，墙上一扇小门}
\BookLine{123.50}{0.00}{\kaifont}{0}{可通往仆人们的住所。}
\BookLine{132.00}{7.38}{\kaifont}{0}{舞台中央前区有一条小路（向右逐渐走低）直接通向剧院}
\BookLine{140.50}{0.00}{\kaifont}{0}{的乐池。这条路的下方是一个“洞穴”——一处有活水流动的}
\BookLine{149.00}{0.00}{\kaifont}{0}{清泉，一扇铜门通往冥界，门被藤条覆盖，却也足够大，能}
\BookLine{157.50}{0.00}{\kaifont}{0}{够让演员从中通过。这儿还有一条（看不到，只存在于想象中}
\BookLine{166.00}{0.00}{\kaifont}{0}{的）蛇——佩索。}
\newpage
% 原书页码：26
\null
\BookNumber{26}
\Vertical{4}{169}{}{6.2}{8}阿\\尔\\刻\\提\\斯\\之\\歌\EndVertical
\BookLine{30.00}{7.38}{\kaifont}{0}{黎明的第一道光线。}
\BookLine{38.50}{7.38}{\kaifont}{0}{一束光渐亮，阿波罗正站在宫殿的房顶上。他穿着金色}
\BookLine{47.00}{0.00}{\kaifont}{0}{的衣服，右肩上披着深蓝色的长斗篷。一束蓝色的光渐渐打}
\BookLine{55.50}{0.00}{\kaifont}{0}{亮泉水下方的冥界入口处，死神——穿着一件黑色的补丁外}
\BookLine{64.00}{0.00}{\kaifont}{0}{套，看上去像一只蝙蝠，或者甲虫——蠕动着沿着小路爬上}
\BookLine{72.50}{0.00}{\kaifont}{0}{来，在左边的门口以及宫殿的门口嗅个不停。}
\BookLine{81.00}{7.38}{\kaifont}{0}{这场戏自始至终，阿波罗一直注视着渐渐升起的太}
\BookLine{89.50}{0.00}{\kaifont}{0}{阳——平静而克制，嘴角带着一丝不易察觉的微笑。}
\BookLine{106.50}{0.00}{\heifont}{0}{阿波罗：}
\BookLine{115.00}{22.14}{}{0}{{\kaifont （像在问早安一样）}死神！}
\BookLine{132.00}{0.00}{\heifont}{0}{死神：}
\BookLine{140.50}{22.14}{}{0}{啊！你来了！阿德墨托斯的宫殿迎来了一位尊}
\BookLine{149.00}{22.14}{}{0}{贵的客人啊！我们这里今天要举行一场婚礼，}
\BookLine{157.50}{22.14}{}{0}{竟有你来做客！你是不是来偷新娘子的啊，大}
\BookLine{166.00}{22.14}{}{0}{名鼎鼎的阿波罗？}
\newpage
% 原书页码：27
\null
\BookNumber{27}
\Vertical{129}{11}{}{7}{10}第\\一\\幕\EndVertical
\BookLine{30.00}{0.00}{\heifont}{0}{阿波罗：}
\BookLine{38.50}{22.14}{}{0}{死神，你活在黑暗之中。}
\BookLine{55.50}{0.00}{\heifont}{0}{死神：}
\BookLine{64.00}{22.14}{}{0}{是啊，是啊。阿波罗大人，你今天来，是不是}
\BookLine{72.50}{22.14}{}{0}{要给我们展现一些伟大的启示，伟大的神迹？}
\BookLine{89.50}{7.38}{\kaifont}{0}{看门人摇着手中的鼓，提着羊皮灯笼绕过宫殿走到舞台}
\BookLine{98.00}{0.00}{\kaifont}{0}{中央。}
\BookLine{115.00}{0.00}{\heifont}{0}{看门人：}
\BookLine{123.50}{22.14}{}{0}{{\kaifont （念词）}黎明时分，照例巡查。热情诚挚的阿德}
\BookLine{132.00}{22.14}{}{0}{墨托斯，塞萨利之王，千金万马的拥有者，宫}
\BookLine{140.50}{22.14}{}{0}{中万事顺遂。{\kaifont （向左侧走下台）}黎明时分，大婚}
\BookLine{149.00}{22.14}{}{0}{之日——将是最伟大的婚礼之一。}
\BookLine{166.00}{7.38}{\kaifont}{0}{他从左侧下场，回到仆人们的住所。}
\newpage
% 原书页码：28
\null
\BookNumber{28}
\Vertical{4}{169}{}{6.2}{8}阿\\尔\\刻\\提\\斯\\之\\歌\EndVertical
\BookLine{30.00}{0.00}{\heifont}{0}{死神：}
\BookLine{38.50}{22.14}{}{0}{我刚问你呢，阿波罗大人，你是不是来给我们}
\BookLine{47.00}{22.14}{}{0}{点奇迹看看的？{\kaifont （停顿）}是？不是？“是。”我希}
\BookLine{55.50}{22.14}{}{0}{望是。当众神靠近人类之时，总是会有人死去。}
\BookLine{64.00}{22.14}{}{0}{我的王国今天是不是会有幸迎来某个尊贵的客}
\BookLine{72.50}{22.14}{}{0}{人？是阿德墨托斯国王还是阿尔刻提斯公主？}
\BookLine{89.50}{0.00}{\heifont}{0}{阿波罗：}
\BookLine{98.00}{22.14}{}{0}{没有。}
\BookLine{115.00}{0.00}{\heifont}{0}{死神：}
\BookLine{123.50}{22.14}{}{0}{让我等着瞧瞧。{\kaifont （他蠕动到场中央）}我能不能}
\BookLine{132.00}{22.14}{}{0}{再问问，你今天是带着哪一种属性和能力出现}
\BookLine{140.50}{22.14}{}{0}{的？治疗者？{\kaifont （停顿）}光明和生命的赐予者？}
\BookLine{149.00}{22.14}{}{0}{{\kaifont （停顿）}歌唱者？}
\BookLine{166.00}{0.00}{\heifont}{0}{阿波罗：}
\BookLine{174.50}{22.14}{}{0}{{\kaifont （仍然注视着太阳，轻松地）}他们是同一个人，}
\newpage
% 原书页码：29
\null
\BookNumber{29}
\Vertical{129}{11}{}{7}{10}第\\一\\幕\EndVertical
\BookLine{30.00}{22.14}{}{0}{是同样的。我今天来，是要开启一曲有情节的}
\BookLine{38.50}{22.14}{}{0}{歌谣——一个故事——}
\BookLine{55.50}{0.00}{\heifont}{0}{死神：}
\BookLine{64.00}{22.14}{}{0}{一个故事！}
\BookLine{81.00}{0.00}{\heifont}{0}{阿波罗：}
\BookLine{89.50}{22.14}{}{0}{一个会被不断传讲的故事。}
\BookLine{106.50}{0.00}{\heifont}{0}{死神：}
\BookLine{115.00}{22.14}{}{0}{啊！那就是一个教导了！其中有给我的教}
\BookLine{123.50}{22.14}{}{0}{导吗？}
\BookLine{140.50}{0.00}{\heifont}{0}{阿波罗：}
\BookLine{149.00}{22.14}{}{0}{是的。}
\BookLine{166.00}{0.00}{\heifont}{0}{死神：}
\BookLine{174.50}{22.14}{}{0}{{\kaifont （用他的手脚拍打地面）}不！}
\newpage
% 原书页码：30
\null
\BookNumber{30}
\Vertical{4}{169}{}{6.2}{8}阿\\尔\\刻\\提\\斯\\之\\歌\EndVertical
\BookLine{30.00}{0.00}{\heifont}{0}{阿波罗：}
\BookLine{38.50}{22.14}{}{0}{是的，你会得到你的教导。}
\BookLine{55.50}{0.00}{\heifont}{0}{死神：}
\BookLine{64.00}{22.14}{}{0}{{\kaifont （勃然大怒）}不！你没资格教导我。我永远都在}
\BookLine{72.50}{22.14}{}{0}{这儿，而且永不改变。你们这些天上的神才需}
\BookLine{81.00}{22.14}{}{0}{要教导。我现在就先给你上一课，{\kaifont （尖声尖气}}
\BookLine{89.50}{22.14}{}{0}{{\kaifont 地）}别管这些人类了，好好待在你的奥林匹斯}
\BookLine{98.00}{22.14}{}{0}{山上，好好过你们的日子。你们的这些蠢念头}
\BookLine{106.50}{22.14}{}{0}{我早就知道了，你们创造了这些生物，然后又}
\BookLine{115.00}{22.14}{}{0}{迷上了他们。结果你们让整个世界乱成了一团，}
\BookLine{123.50}{22.14}{}{0}{一天比一天糟糕。你所做的对他们不过是种折}
\BookLine{132.00}{22.14}{}{0}{磨——谁能比我更了解呢？{\kaifont （他往回蠕动到小}}
\BookLine{140.50}{22.14}{}{0}{{\kaifont 路的最上方，狂暴地晃动着身体）}他们永远不}
\BookLine{149.00}{22.14}{}{0}{会理解你的意思。你越是想要表达，就会把他}
\BookLine{157.50}{22.14}{}{0}{们弄得越疯狂。}
\newpage
% 原书页码：31
\null
\BookNumber{31}
\Vertical{129}{11}{}{7}{10}第\\一\\幕\EndVertical
\BookLine{30.00}{0.00}{\heifont}{0}{阿波罗：}
\BookLine{38.50}{22.14}{}{0}{他们已经开始理解我了。最初他们的确就像野}
\BookLine{47.00}{22.14}{}{0}{兽——野蛮，恐惧。就像是笼中的困兽，他们}
\BookLine{55.50}{22.14}{}{0}{自己就是自己的牢笼。后来有两件事情让他们}
\BookLine{64.00}{22.14}{}{0}{开始有了意识，开始寻索：我父亲的雷电，让}
\BookLine{72.50}{22.14}{}{0}{他们的恐惧变为敬畏；我的阳光，让他们学会}
\BookLine{81.00}{22.14}{}{0}{了感谢。在感谢之中他们发明了语言，我又给}
\BookLine{89.50}{22.14}{}{0}{了他们音乐。这些就是启示，他们理解了。一}
\BookLine{98.00}{22.14}{}{0}{个又一个人，意识到是我在搅动他们的心灵，}
\BookLine{106.50}{22.14}{}{0}{对他们言语。}
\BookLine{123.50}{0.00}{\heifont}{0}{死神：}
\BookLine{132.00}{22.14}{}{0}{是的，他们和以前不一样了：“阿波罗爱塞萨}
\BookLine{140.50}{22.14}{}{0}{利，阿波罗爱费莱。”回你的奥林匹斯山去吧。}
\BookLine{149.00}{22.14}{}{0}{这种爱吧……很难说你们谁更不幸——是你，}
\BookLine{157.50}{22.14}{}{0}{还是这些可怜的造物。你想要介入他们的生活，}
\BookLine{166.00}{22.14}{}{0}{这就像一个巨人走进了小小的房间：你每挪一}
\BookLine{174.50}{22.14}{}{0}{小步，都会弄坏点什么。你今天打算让谁吃这}
\newpage
% 原书页码：32
\null
\BookNumber{32}
\Vertical{4}{169}{}{6.2}{8}阿\\尔\\刻\\提\\斯\\之\\歌\EndVertical
\BookLine{30.00}{22.14}{}{0}{个苦头呢？国王？还是他的新娘？}
\BookLine{47.00}{0.00}{\heifont}{0}{阿波罗：}
\BookLine{55.50}{22.14}{}{0}{你。}
\BookLine{72.50}{0.00}{\heifont}{0}{死神：}
\BookLine{81.00}{22.14}{}{0}{我？我？所以你是打算也爱一爱我喽？不！谢}
\BookLine{89.50}{22.14}{}{0}{谢！{\kaifont （移动起所有手脚，在小路上来回爬着，尖}}
\BookLine{98.00}{22.14}{}{0}{{\kaifont 叫着）}你休想找我麻烦，你也休想教导我。我}
\BookLine{106.50}{22.14}{}{0}{和我的王国会永远存在。你怎么可能给我找麻}
\BookLine{115.00}{22.14}{}{0}{烦呢？}
\BookLine{132.00}{0.00}{\heifont}{0}{阿波罗：}
\BookLine{140.50}{22.14}{}{0}{你生活在黑暗之中，所以看不到万事万物的}
\BookLine{149.00}{22.14}{}{0}{改变。}
\BookLine{166.00}{0.00}{\heifont}{0}{死神：}
\BookLine{174.50}{22.14}{}{0}{{\kaifont （喊叫着）}改变！没什么会变！{\kaifont （恐惧地环视一}}
\newpage
% 原书页码：33
\null
\BookNumber{33}
\Vertical{129}{11}{}{7}{10}第\\一\\幕\EndVertical
\BookLine{30.00}{22.14}{}{0}{{\kaifont 下四周）}天马上要亮了。所以你今天要讲的这}
\BookLine{38.50}{22.14}{}{0}{个故事是关于改变？针对我的改变？}
\BookLine{55.50}{0.00}{\heifont}{0}{阿波罗：}
\BookLine{64.00}{22.14}{}{0}{对你，也对我。}
\BookLine{81.00}{0.00}{\heifont}{0}{死神：}
\BookLine{89.50}{22.14}{}{0}{{\kaifont （消失在自己的洞穴，最后尖叫道）}就对你！}
\BookLine{106.50}{7.38}{\kaifont}{0}{阿波罗身上的光渐渐隐去，他消失在大家的视野中。}
\BookLine{115.00}{7.38}{\kaifont}{0}{看门人回到场上，他摇了摇鼓，把灯吹灭。}
\BookLine{132.00}{0.00}{\heifont}{0}{看门人：}
\BookLine{140.50}{22.14}{}{0}{黎明到来，黎明到来。热情诚挚的阿德墨托}
\BookLine{149.00}{22.14}{}{0}{斯，塞萨利之王，千金万马的拥有者，宫中万}
\BookLine{157.50}{22.14}{}{0}{事顺遂。{\kaifont （走向观众）}大婚之日——将是最伟}
\BookLine{166.00}{22.14}{}{0}{大的婚礼之一，可事情都不顺利。她为什么无}
\BookLine{174.50}{22.14}{}{0}{法入睡——这位公主，这位新娘，我们未来的}
\newpage
% 原书页码：34
\null
\BookNumber{34}
\Vertical{4}{169}{}{6.2}{8}阿\\尔\\刻\\提\\斯\\之\\歌\EndVertical
\BookLine{30.00}{22.14}{}{0}{王后？就在夜里，我曾八次、十次在这里见到}
\BookLine{38.50}{22.14}{}{0}{她——望着天空四处游荡。有时她走到大路上}
\BookLine{47.00}{22.14}{}{0}{去，好像在等着什么信使。她站在这里，举起}
\BookLine{55.50}{22.14}{}{0}{她的双臂低语：“阿波罗！一点启示！一个启}
\BookLine{64.00}{22.14}{}{0}{示！”启示什么？她嫁给阿德墨托斯国王是否}
\BookLine{72.50}{22.14}{}{0}{正确？唉！她想寻求的最清楚明白的启示早已}
\BookLine{81.00}{22.14}{}{0}{写在国王脸上了。哦，我已经有些岁数了，我}
\BookLine{89.50}{22.14}{}{0}{知道每一位新娘在她们的婚礼前夜内心都会充}
\BookLine{98.00}{22.14}{}{0}{满不安，但是何必为我们的阿德墨托斯国王感}
\BookLine{106.50}{22.14}{}{0}{到不安，他赢得她的奇妙故事让整个希腊都为}
\BookLine{115.00}{22.14}{}{0}{之惊讶。哦，朋友们，请记住我这个老看门人}
\BookLine{123.50}{22.14}{}{0}{的建议：不要在凌晨三点的时候思考人生大事。}
\BookLine{132.00}{22.14}{}{0}{在那样的时刻，你的心灵和头脑都没有温度，}
\BookLine{140.50}{22.14}{}{0}{在那样的时刻——嘘——你会看到自己的房子}
\BookLine{149.00}{22.14}{}{0}{着起火，你的孩子挣扎着死在你的脚边。等到}
\BookLine{157.50}{22.14}{}{0}{太阳升起吧，事实都是一样的——人类生活的}
\BookLine{166.00}{22.14}{}{0}{一切事实都是一样的——但阳光赋予了它们意}
\BookLine{174.50}{22.14}{}{0}{义。请记住看门人给你们的建议。现在我要去}
\newpage
% 原书页码：35
\null
\BookNumber{35}
\Vertical{129}{11}{}{7}{10}第\\一\\幕\EndVertical
\BookLine{30.00}{22.14}{}{0}{喝点水。{\kaifont （走到泉水旁边，和蛇佩索打了个招}}
\BookLine{38.50}{22.14}{}{0}{{\kaifont 呼）}日安佩索，我的老朋友。今天你一定会过}
\BookLine{47.00}{22.14}{}{0}{得不错，婚宴少不了你的份儿——一只羊，或}
\BookLine{55.50}{22.14}{}{0}{者是半头牛。过去，好让我进行我的祭礼。{\kaifont （让}}
\BookLine{64.00}{22.14}{}{0}{{\kaifont 手中捧着的水流出，一边念着祷文：“生命之源}}
\BookLine{72.50}{22.14}{}{0}{{\kaifont 啊——大地，空气，火，水……”然后又捧了}}
\BookLine{81.00}{22.14}{}{0}{{\kaifont 一捧水喝了下去。对观众说话）}看，我的朋友}
\BookLine{89.50}{22.14}{}{0}{们。你们看到这些藤条下面的洞穴了吗？这是}
\BookLine{98.00}{22.14}{}{0}{前往冥界的五个入口之一，我们的好佩索正是}
\BookLine{106.50}{22.14}{}{0}{在此守护着入口。没有人进到过那里，也没有}
\BookLine{115.00}{22.14}{}{0}{人从那里出来。就是这样——我们有很多不明}
\BookLine{123.50}{22.14}{}{0}{白的事，这只是万分之一。{\kaifont （又喝了点水，同}}
\BookLine{132.00}{22.14}{}{0}{{\kaifont 时后面传来动物被宰杀时的尖叫）}好啦，大日}
\BookLine{140.50}{22.14}{}{0}{子开始了。他们已经在为宴会宰杀牲畜了。厨}
\BookLine{149.00}{22.14}{}{0}{子们生起了火，草原上满是国王们和首领们的}
\BookLine{157.50}{22.14}{}{0}{帐篷，他们都是来庆祝阿德墨托斯国王和伊俄}
\BookLine{166.00}{22.14}{}{0}{尔科斯国王佩利阿斯之女阿尔刻提斯公主的婚}
\BookLine{174.50}{22.14}{}{0}{礼的。}
\newpage
% 原书页码：36
\null
\BookNumber{36}
\Vertical{4}{169}{}{6.2}{8}阿\\尔\\刻\\提\\斯\\之\\歌\EndVertical
\BookLine{30.00}{7.38}{\kaifont}{0}{他沿着小路往回走着。阿尔刻提斯穿着白色衣服从宫门}
\BookLine{38.50}{0.00}{\kaifont}{0}{缓步走出。可以听到动物们叫唤的声音。}
\BookLine{55.50}{0.00}{\heifont}{0}{看门人：}
\BookLine{64.00}{22.14}{}{0}{{\kaifont （继续说道）}嘘，她又来了！}
\BookLine{81.00}{7.38}{\kaifont}{0}{看门人在路边略低于舞台的位置藏了起来。}
\BookLine{89.50}{7.38}{\kaifont}{0}{阿尔刻提斯走到舞台中央，举起双臂，低语着。}
\BookLine{106.50}{0.00}{\heifont}{0}{阿尔刻提斯：}
\BookLine{115.00}{22.14}{}{0}{阿波罗！一点启示！一个启示！{\kaifont （走出大门，向}}
\BookLine{123.50}{22.14}{}{0}{{\kaifont 右朝大路上走去）}}
\BookLine{140.50}{0.00}{\heifont}{0}{看门人：}
\BookLine{149.00}{22.14}{}{0}{{\kaifont （模仿她，对着观众轻轻地喊着）}阿波罗！一点}
\BookLine{157.50}{22.14}{}{0}{启示！}
\BookLine{174.50}{7.38}{\kaifont}{0}{动物们仍然在叫唤。老奶妈阿格莱亚从宫殿里跑出来。}
\newpage
% 原书页码：37
\null
\BookNumber{37}
\Vertical{129}{11}{}{7}{10}第\\一\\幕\EndVertical
\BookLine{30.00}{0.00}{\kaifont}{0}{她四处张望，看见了看门人。}
\BookLine{47.00}{0.00}{\heifont}{0}{阿格莱亚：}
\BookLine{55.50}{22.14}{}{1}{{\kaifont （小声地）}公主在哪儿？{\kaifont （看向看门人指的方向）}}
\BookLine{64.00}{22.14}{}{0}{一整个晚上——一整晚不眠不休，忧伤满怀！}
\BookLine{72.50}{22.14}{}{0}{“阿波罗！一点启示！”}
\BookLine{89.50}{0.00}{\heifont}{0}{看门人：}
\BookLine{98.00}{22.14}{}{0}{“阿波罗！一点启示！”}
\BookLine{115.00}{7.38}{\kaifont}{0}{阿尔刻提斯从右侧的门走进来，带着一种神经质的果决}
\BookLine{123.50}{0.00}{\kaifont}{0}{情绪。}
\BookLine{140.50}{0.00}{\heifont}{0}{阿尔刻提斯：}
\BookLine{149.00}{22.14}{}{0}{看门人！}
\BookLine{166.00}{0.00}{\heifont}{0}{看门人：}
\BookLine{174.50}{22.14}{}{0}{在，公主？}
\newpage
% 原书页码：38
\null
\BookNumber{38}
\Vertical{4}{169}{}{6.2}{8}阿\\尔\\刻\\提\\斯\\之\\歌\EndVertical
\BookLine{30.00}{0.00}{\heifont}{0}{阿尔刻提斯：}
\BookLine{38.50}{22.14}{}{0}{把我的马夫都找来。吩咐他们把马都备好，我}
\BookLine{47.00}{22.14}{}{0}{要出门。阿格莱亚，把我的女仆都找来，吩咐}
\BookLine{55.50}{22.14}{}{0}{她们把行李都准备好。}
\BookLine{72.50}{0.00}{\heifont}{0}{阿格莱亚：}
\BookLine{81.00}{22.14}{}{0}{您要出门，公主——在您大婚的日子，要}
\BookLine{89.50}{22.14}{}{0}{出门？}
\BookLine{106.50}{0.00}{\heifont}{0}{阿尔刻提斯：}
\BookLine{115.00}{22.14}{}{0}{{\kaifont （迅速走下台阶到她面前）}阿格莱亚，我别无选}
\BookLine{123.50}{22.14}{}{0}{择。我必须离开。原谅我。不，恨我吧，蔑视}
\BookLine{132.00}{22.14}{}{0}{我吧——不过最终，请忘了我。}
\BookLine{149.00}{0.00}{\heifont}{0}{阿格莱亚：}
\BookLine{157.50}{22.14}{}{0}{公主，这是个耻辱，这对阿德墨托斯国王是种}
\BookLine{166.00}{22.14}{}{0}{侮辱。}
\newpage
% 原书页码：39
\null
\BookNumber{39}
\Vertical{129}{11}{}{7}{10}第\\一\\幕\EndVertical
\BookLine{30.00}{0.00}{\heifont}{0}{阿尔刻提斯：}
\BookLine{38.50}{22.14}{}{0}{我都明白，阿格莱亚。阿格莱亚，等我离开之}
\BookLine{47.00}{22.14}{}{0}{后，告诉国王——告诉阿德墨托斯——所有的}
\BookLine{55.50}{22.14}{}{0}{耻辱我独自承担，我不奢求他可以原谅，让他}
\BookLine{64.00}{22.14}{}{0}{蔑视我，然后忘了我吧。}
\BookLine{81.00}{0.00}{\heifont}{0}{阿格莱亚：}
\BookLine{89.50}{22.14}{}{0}{公主，我是个老妇人了，但我不是宫里一个普}
\BookLine{98.00}{22.14}{}{0}{通的仆人。我养育了幼年的阿德墨托斯，还有}
\BookLine{106.50}{22.14}{}{0}{他的父亲。{\kaifont （对看门人）}看门人，请留我们单独}
\BookLine{115.00}{22.14}{}{0}{说话。}
\BookLine{132.00}{7.38}{\kaifont}{0}{看门人下场。}
\BookLine{149.00}{0.00}{\heifont}{0}{阿格莱亚：}
\BookLine{157.50}{22.14}{}{0}{您不了解阿德墨托斯国王。希腊众岛之中，您}
\BookLine{166.00}{22.14}{}{0}{再找不到比他更好的丈夫了。}
\newpage
% 原书页码：40
\null
\BookNumber{40}
\Vertical{4}{169}{}{6.2}{8}阿\\尔\\刻\\提\\斯\\之\\歌\EndVertical
\BookLine{30.00}{0.00}{\heifont}{0}{阿尔刻提斯：}
\BookLine{38.50}{22.14}{}{0}{我知道。}
\BookLine{55.50}{0.00}{\heifont}{0}{阿格莱亚：}
\BookLine{64.00}{22.14}{}{0}{或许有人比他更善战、更进取、更强壮，但不}
\BookLine{72.50}{22.14}{}{0}{会有人比他更——更受人爱戴了。}
\BookLine{89.50}{0.00}{\heifont}{0}{阿尔刻提斯：}
\BookLine{98.00}{22.14}{}{0}{我都知道，阿格莱亚。我，我也爱阿德墨托斯。}
\BookLine{106.50}{22.14}{}{0}{所以我也不知道为什么高兴不起来。可是比起}
\BookLine{115.00}{22.14}{}{0}{他，我更爱另一个。}
\BookLine{132.00}{0.00}{\heifont}{0}{阿格莱亚：}
\BookLine{140.50}{22.14}{}{0}{另一个？另一个男人？比阿德墨托斯更爱？那}
\BookLine{149.00}{22.14}{}{0}{去吧，公主，快去吧。我们错待您了，您不必}
\BookLine{157.50}{22.14}{}{0}{待在这里。既然您有眼无珠，麻木无知，什么}
\BookLine{166.00}{22.14}{}{0}{都看不到。{\kaifont （严厉地）}看门人！看门人！您路上}
\BookLine{174.50}{22.14}{}{0}{所需要的一切都会准备好，公主，快点走吧。}
\newpage
% 原书页码：41
\null
\BookNumber{41}
\Vertical{129}{11}{}{7}{10}第\\一\\幕\EndVertical
\BookLine{30.00}{0.00}{\heifont}{0}{阿尔刻提斯：}
\BookLine{38.50}{22.14}{}{0}{不，阿格莱亚，不是别的男人。比起阿德墨托}
\BookLine{47.00}{22.14}{}{0}{斯，我更爱的那一位是，是一位神祇，是阿}
\BookLine{55.50}{22.14}{}{0}{波罗。}
\BookLine{72.50}{0.00}{\heifont}{0}{阿格莱亚：}
\BookLine{81.00}{22.14}{}{0}{阿波罗？}
\BookLine{98.00}{0.00}{\heifont}{0}{阿尔刻提斯：}
\BookLine{106.50}{22.14}{}{0}{是的。从我还是个小姑娘起，我就只有一个心}
\BookLine{115.00}{22.14}{}{0}{愿——在德尔斐，做他的女祭司。{\kaifont （越说越沮}}
\BookLine{123.50}{22.14}{}{0}{{\kaifont 丧，突然哭了起来，很激动地）}我想过真正的}
\BookLine{132.00}{22.14}{}{0}{生活。只此一生，奉献所有。——不是我们所}
\BookLine{140.50}{22.14}{}{0}{看到的狂热、骄傲、占有的生活，而是真正的}
\BookLine{149.00}{22.14}{}{0}{生活，在德尔斐，在真理所在之处。}
\BookLine{166.00}{0.00}{\heifont}{0}{阿格莱亚：}
\BookLine{174.50}{22.14}{}{0}{可是，神呼召您了吗？{\kaifont （停顿）}神没有差遣过谁}
\newpage
% 原书页码：42
\null
\BookNumber{42}
\Vertical{4}{169}{}{6.2}{8}阿\\尔\\刻\\提\\斯\\之\\歌\EndVertical
\BookLine{30.00}{22.14}{}{0}{来找您吗？}
\BookLine{47.00}{0.00}{\heifont}{0}{阿尔刻提斯：}
\BookLine{55.50}{22.14}{}{0}{{\kaifont （低声，羞愧地）}没有。}
\BookLine{72.50}{0.00}{\heifont}{0}{阿格莱亚：}
\BookLine{81.00}{22.14}{}{0}{所谓真正的——难道成为阿德墨托斯的妻子、}
\BookLine{89.50}{22.14}{}{0}{孩子们的母亲、塞萨利的王后，就不算真正的}
\BookLine{98.00}{22.14}{}{0}{生活吗？}
\BookLine{115.00}{0.00}{\heifont}{0}{阿尔刻提斯：}
\BookLine{123.50}{22.14}{}{0}{每个女人都可以做妻子，做母亲，成百上千人}
\BookLine{132.00}{22.14}{}{0}{做过王后。然后，我的丈夫，我的孩子……生}
\BookLine{140.50}{22.14}{}{0}{活就开始围绕这五六个人打转，我只会在意与}
\BookLine{149.00}{22.14}{}{0}{他们相关的琐事……每一天就这样过去了——}
\BookLine{157.50}{22.14}{}{0}{被塞满——被他们无数的需求塞满，最终带}
\BookLine{166.00}{22.14}{}{0}{着点爱和尊荣躺进坟墓——仍然和出生时一样}
\BookLine{174.50}{22.14}{}{0}{无知——}
\newpage
% 原书页码：43
\null
\BookNumber{43}
\Vertical{129}{11}{}{7}{10}第\\一\\幕\EndVertical
\BookLine{30.00}{0.00}{\heifont}{0}{阿格莱亚：}
\BookLine{38.50}{22.14}{}{0}{无知？}
\BookLine{55.50}{0.00}{\heifont}{0}{阿尔刻提斯：}
\BookLine{64.00}{22.14}{}{0}{对我们为何出生，为何死去——人们为何出生}
\BookLine{72.50}{22.14}{}{0}{死去——一无所知，和我们出生时一样无知。}
\BookLine{89.50}{0.00}{\heifont}{0}{阿格莱亚：}
\BookLine{98.00}{22.14}{}{0}{{\kaifont （不带感情地）}所以，您认为在德尔斐可以知}
\BookLine{106.50}{22.14}{}{0}{道？但神并没有呼召您啊。}
\BookLine{123.50}{0.00}{\heifont}{0}{阿尔刻提斯：}
\BookLine{132.00}{22.14}{}{0}{{\kaifont （羞愧地）}我送去了奉献，带去了信息，献了}
\BookLine{140.50}{22.14}{}{0}{祭……{\kaifont （停顿）}我是父亲最疼爱的女儿，他曾}
\BookLine{149.00}{22.14}{}{0}{希望我终身不嫁，一直陪伴他左右。但追求者}
\BookLine{157.50}{22.14}{}{0}{们从希腊各个地方涌来，于是父亲提出了一项}
\BookLine{166.00}{22.14}{}{0}{不可能完成的任务。他让他们同时拴住一头狮}
\BookLine{174.50}{22.14}{}{0}{子和一只野猪，并驾着它们沿伊俄尔科斯城的}
\newpage
% 原书页码：44
\null
\BookNumber{44}
\Vertical{4}{169}{}{6.2}{8}阿\\尔\\刻\\提\\斯\\之\\歌\EndVertical
\BookLine{30.00}{22.14}{}{0}{城墙绕满三周。人们从希腊各个地方涌来：伊}
\BookLine{38.50}{22.14}{}{0}{阿宋来了，涅斯托耳来了，宙斯之子赫拉克勒}
\BookLine{47.00}{22.14}{}{0}{斯也来了，还有阿特柔斯。可所有人都失败了。}
\BookLine{55.50}{22.14}{}{0}{几个月过去，不断有新的追求者失败甚至丧命}
\BookLine{64.00}{22.14}{}{0}{于此。父亲和我就坐在城门之上，父亲笑着，}
\BookLine{72.50}{22.14}{}{0}{我也笑着。不是因为我渴望留在父亲身边，而}
\BookLine{81.00}{22.14}{}{0}{是因为我所求的那唯一一件事——此生只在德}
\BookLine{89.50}{22.14}{}{0}{尔斐，做阿波罗的女祭司。}
\BookLine{106.50}{0.00}{\heifont}{0}{阿格莱亚：}
\BookLine{115.00}{22.14}{}{0}{然后阿德墨托斯出现了。他驾驭着狮子和野}
\BookLine{123.50}{22.14}{}{0}{猪——就像驾驭着温顺的公牛一样绕城而行，}
\BookLine{132.00}{22.14}{}{0}{并赢得了您，公主。}
\BookLine{149.00}{0.00}{\heifont}{0}{阿尔刻提斯：}
\BookLine{157.50}{22.14}{}{0}{可我更爱阿波罗。}
\newpage
% 原书页码：45
\null
\BookNumber{45}
\Vertical{129}{11}{}{7}{10}第\\一\\幕\EndVertical
\BookLine{30.00}{0.00}{\heifont}{0}{阿格莱亚：}
\BookLine{38.50}{22.14}{}{0}{是，可正是阿波罗安排了这场婚礼。}
\BookLine{55.50}{0.00}{\heifont}{0}{阿尔刻提斯：}
\BookLine{64.00}{22.14}{}{0}{那可说不准。}
\BookLine{81.00}{0.00}{\heifont}{0}{阿格莱亚：}
\BookLine{89.50}{22.14}{}{0}{您要求的启示，公主，就在您的面前——最清}
\BookLine{98.00}{22.14}{}{0}{楚无误的启示。{\kaifont （激烈地）}您从没有被呼召前往}
\BookLine{106.50}{22.14}{}{0}{德尔斐，难道您看不出神最简单的指引吗？}
\BookLine{123.50}{7.38}{\kaifont}{0}{阿尔刻提斯捧起了自己的面庞。}
\BookLine{140.50}{0.00}{\heifont}{0}{阿格莱亚：}
\BookLine{149.00}{22.14}{}{0}{现在让我告诉您：阿德墨托斯能够驾驭狮子与}
\BookLine{157.50}{22.14}{}{0}{野猪，难道您不惊讶吗？阿特柔斯失败了，宙}
\BookLine{166.00}{22.14}{}{0}{斯之子赫拉克勒斯也失败了，他却做到了。我}
\newpage
% 原书页码：46
\null
\BookNumber{46}
\Vertical{4}{169}{}{6.2}{8}阿\\尔\\刻\\提\\斯\\之\\歌\EndVertical
\BookLine{30.00}{22.14}{}{0}{来告诉您他是怎么做到的：阿波罗曾在梦中教}
\BookLine{38.50}{22.14}{}{0}{授他如何同时驾驭一只狮子和一只野猪。}
\BookLine{55.50}{7.38}{\kaifont}{0}{阿尔刻提斯向后退了两步。}
\BookLine{72.50}{0.00}{\heifont}{0}{阿格莱亚：}
\BookLine{81.00}{22.14}{}{0}{起初，他对您一见钟情。他第二次回到伊俄尔}
\BookLine{89.50}{22.14}{}{0}{科斯之前——也就是第一次失败后，他病了。}
\BookLine{98.00}{22.14}{}{0}{爱情与绝望让他濒临死亡，我一直在身旁照顾}
\BookLine{106.50}{22.14}{}{0}{他，有三个晚上他都差点死去。就在第三晚，}
\BookLine{115.00}{22.14}{}{0}{我就坐在他身边——在痛苦与谵妄之中——我}
\BookLine{123.50}{22.14}{}{0}{听到，我看到阿波罗正在梦中教授他，如何同}
\BookLine{132.00}{22.14}{}{0}{时驾驭一只狮子和一只野猪。}
\BookLine{149.00}{7.38}{\kaifont}{0}{阿尔刻提斯盯着她。}
\BookLine{166.00}{0.00}{\heifont}{0}{阿格莱亚：}
\BookLine{174.50}{22.14}{}{0}{这是真的。我发誓这都是真的。}
\newpage
% 原书页码：47
\null
\BookNumber{47}
\Vertical{129}{11}{}{7}{10}第\\一\\幕\EndVertical
\BookLine{30.00}{0.00}{\heifont}{0}{阿尔刻提斯：}
\BookLine{38.50}{22.14}{}{0}{是真的，好的。不过我们听说过许多关于谵妄}
\BookLine{47.00}{22.14}{}{0}{和梦境、发热和幻象的事情。阿格莱亚，我们}
\BookLine{55.50}{22.14}{}{0}{必须百分百确定。神清晰地公然降临——只会}
\BookLine{64.00}{22.14}{}{0}{是在德尔斐。}
\BookLine{81.00}{0.00}{\heifont}{0}{阿格莱亚：}
\BookLine{89.50}{22.14}{}{0}{清晰？公然？即便在德尔斐，女预言家们也都}
\BookLine{98.00}{22.14}{}{0}{处在谵妄之中，她们高声咆哮，魂游象外。谁}
\BookLine{106.50}{22.14}{}{0}{听到过她们清晰明白地讲话？}
\BookLine{123.50}{0.00}{\heifont}{0}{阿尔刻提斯：}
\BookLine{132.00}{22.14}{}{0}{{\kaifont （拖着犹豫不决的步子朝王宫走去，十分沮丧）}}
\BookLine{140.50}{22.14}{}{0}{我孤独一人，只有一人……}
\BookLine{157.50}{0.00}{\heifont}{0}{阿格莱亚：}
\BookLine{166.00}{22.14}{}{0}{{\kaifont （坚定地，动情地）}听我说，公主。回到房间好}
\BookLine{174.50}{22.14}{}{0}{好睡一觉。{\kaifont （朝天看一眼）}还有两个钟头就到正}
\newpage
% 原书页码：48
\null
\BookNumber{48}
\Vertical{4}{169}{}{6.2}{8}阿\\尔\\刻\\提\\斯\\之\\歌\EndVertical
\BookLine{30.00}{22.14}{}{0}{午了。如果休息之后，您的想法还是没有改变，}
\BookLine{38.50}{22.14}{}{0}{您想去哪里就去哪里，没有人会阻拦。{\kaifont （托起阿}}
\BookLine{47.00}{22.14}{}{0}{{\kaifont 尔刻提斯的胳膊，把她领进宫门，像母亲和女}}
\BookLine{55.50}{22.14}{}{0}{{\kaifont 儿闲聊一样）}您希望神明们清晰坦率地对我们}
\BookLine{64.00}{22.14}{}{0}{说话，公主？您把他们当作什么了啊，公主？}
\BookLine{72.50}{22.14}{}{0}{您可不能把他们想象成普通的人啊！}
\BookLine{89.50}{7.38}{\kaifont}{0}{两人都退回宫殿中。舞台上空了一段时间。门外一阵声}
\BookLine{98.00}{0.00}{\kaifont}{0}{响愈加吵闹起来，有人在敲右侧的门。看门人绕过宫殿走到}
\BookLine{106.50}{0.00}{\kaifont}{0}{舞台中央。}
\BookLine{123.50}{0.00}{\heifont}{0}{看门人：}
\BookLine{132.00}{22.14}{}{0}{好吧，又是什么情况啊？吵吵嚷嚷是怎么回}
\BookLine{140.50}{22.14}{}{0}{事？{\kaifont （开门，透过缝隙交谈）}婚礼的客人们请走}
\BookLine{149.00}{22.14}{}{0}{宫殿正门。这是后门。什么？能不能不要同时}
\BookLine{157.50}{22.14}{}{0}{讲话！什么？好吧，好吧，让老人家进来。}
\BookLine{174.50}{7.38}{\kaifont}{0}{忒瑞西阿斯走了进来，他是个瞎子，非常非常老，容易}
\newpage
% 原书页码：49
\null
\BookNumber{49}
\Vertical{129}{11}{}{7}{10}第\\一\\幕\EndVertical
\BookLine{30.00}{0.00}{\kaifont}{0}{动怒，言辞恶毒，盛气凌人，因为高龄，讲起话像是胡言乱}
\BookLine{38.50}{0.00}{\kaifont}{0}{语。他一只手搭在带领着自己的男孩肩上，另一只手则不停}
\BookLine{47.00}{0.00}{\kaifont}{0}{挥舞着手中的手杖。镇上的人们都跟着他走进院子，仆人们}
\BookLine{55.50}{0.00}{\kaifont}{0}{也从宫门和皇宫四周来到院子里。}
\BookLine{72.50}{0.00}{\heifont}{0}{忒瑞西阿斯：}
\BookLine{81.00}{22.14}{}{0}{{\kaifont （声音倒是令人惊讶地响亮）}这是克里特国王米}
\BookLine{89.50}{22.14}{}{0}{诺斯的王宫吗？}
\BookLine{106.50}{7.38}{\kaifont}{0}{大家笑起来，男孩拽了拽他的袖子并对他耳语几句。}
\BookLine{123.50}{0.00}{\heifont}{0}{忒瑞西阿斯：}
\BookLine{132.00}{22.14}{}{0}{我是说，这是忒拜国王俄狄浦斯的王宫吗？{\kaifont （用}}
\BookLine{140.50}{22.14}{}{0}{{\kaifont 手杖敲打男孩）}别推搡我！我知道自己在说什}
\BookLine{149.00}{22.14}{}{0}{么！{\kaifont （躲开男孩的推搡）}蜜蜂！胡蜂！大黄蜂！}
\BookLine{166.00}{0.00}{\heifont}{0}{看门人：}
\BookLine{174.50}{22.14}{}{0}{老人家。这是塞萨利国王，阿德墨托斯的王宫。}
\newpage
% 原书页码：50
\null
\BookNumber{50}
\Vertical{4}{169}{}{6.2}{8}阿\\尔\\刻\\提\\斯\\之\\歌\EndVertical
\BookLine{30.00}{0.00}{\heifont}{0}{忒瑞西阿斯：}
\BookLine{38.50}{22.14}{}{0}{{\kaifont （重复着他的话）}塞萨利国王。对，我是这么说}
\BookLine{47.00}{22.14}{}{0}{的。我就是这个意思。把塞萨利国王阿德墨托}
\BookLine{55.50}{22.14}{}{0}{斯找来，我有个消息带给他。}
\BookLine{72.50}{0.00}{\heifont}{0}{看门人：}
\BookLine{81.00}{22.14}{}{0}{老人家，国王今天要举行婚礼，他此刻正忙着}
\BookLine{89.50}{22.14}{}{0}{招待客人。您先坐在这里晒晒太阳，我们会给}
\BookLine{98.00}{22.14}{}{0}{您洗洗脚。稍后国王会来听您的消息。}
\BookLine{115.00}{0.00}{\heifont}{0}{忒瑞西阿斯：}
\BookLine{123.50}{22.14}{}{0}{{\kaifont （举着手杖威胁道）}结婚……洗脚……结婚和}
\BookLine{132.00}{22.14}{}{0}{洗脚跟我有什么关系？我一分钟都不会多等。}
\BookLine{140.50}{22.14}{}{0}{{\kaifont （跺脚）}把那个什么什么国王给我找来！}
\BookLine{157.50}{7.38}{\kaifont}{0}{阿格莱亚从宫中走出来。}
\newpage
% 原书页码：51
\null
\BookNumber{51}
\Vertical{129}{11}{}{7}{10}第\\一\\幕\EndVertical
\BookLine{30.00}{0.00}{\heifont}{0}{阿格莱亚：}
\BookLine{38.50}{22.14}{}{0}{老人家，请问您是？我去通报国王。}
\BookLine{55.50}{0.00}{\heifont}{0}{忒瑞西阿斯：}
\BookLine{64.00}{22.14}{}{0}{告诉国王，我是德尔斐，忒瑞西阿斯的祭}
\BookLine{72.50}{22.14}{}{0}{司——我是阿波罗，德尔斐的祭司……孩子，}
\BookLine{81.00}{22.14}{}{0}{我在说什么？}
\BookLine{98.00}{7.38}{\kaifont}{0}{男孩耳语几句。}
\BookLine{115.00}{0.00}{\heifont}{0}{忒瑞西阿斯：}
\BookLine{123.50}{22.14}{}{0}{告诉国王我是忒瑞西阿斯，阿波罗的祭司。我}
\BookLine{132.00}{22.14}{}{0}{从德尔斐来，带来了一个消息。我今天还要赶}
\BookLine{140.50}{22.14}{}{0}{回去。}
\BookLine{157.50}{0.00}{\heifont}{0}{阿格莱亚：}
\BookLine{166.00}{22.14}{}{0}{你是忒瑞西阿斯？忒瑞西阿斯！}
\newpage
% 原书页码：52
\null
\BookNumber{52}
\Vertical{4}{169}{}{6.2}{8}阿\\尔\\刻\\提\\斯\\之\\歌\EndVertical
\BookLine{30.00}{0.00}{\heifont}{0}{看门人和众人：}
\BookLine{38.50}{22.14}{}{0}{忒瑞西阿斯！}
\BookLine{55.50}{0.00}{\heifont}{0}{忒瑞西阿斯：}
\BookLine{64.00}{22.14}{}{0}{{\kaifont （拿手杖敲打地面）}去把国王叫来！灾难和瘟疫}
\BookLine{72.50}{22.14}{}{0}{要来了！去把那个什么什么国王叫来！}
\BookLine{89.50}{0.00}{\heifont}{0}{阿格莱亚：}
\BookLine{98.00}{22.14}{}{0}{这边请，伟大的忒瑞西阿斯。}
\BookLine{115.00}{7.38}{\kaifont}{0}{她匆匆走进宫殿，门开了，阿德墨托斯走了出来，更多}
\BookLine{123.50}{0.00}{\kaifont}{0}{的仆人围了过来。}
\BookLine{140.50}{0.00}{\heifont}{0}{阿德墨托斯：}
\BookLine{149.00}{22.14}{}{0}{出什么事了？阿格莱亚，这是谁？}
\BookLine{166.00}{0.00}{\heifont}{0}{阿格莱亚：}
\BookLine{174.50}{22.14}{}{0}{{\kaifont （悄悄地）}是从德尔斐来的，忒瑞西阿斯。}
\newpage
% 原书页码：53
\null
\BookNumber{53}
\Vertical{129}{11}{}{7}{10}第\\一\\幕\EndVertical
\BookLine{30.00}{0.00}{\heifont}{0}{阿德墨托斯：}
\BookLine{38.50}{22.14}{}{0}{忒瑞西阿斯！}
\BookLine{55.50}{0.00}{\heifont}{0}{阿格莱亚：}
\BookLine{64.00}{22.14}{}{0}{{\kaifont （手指前额）}如山峰一样古老，阿德墨托斯}
\BookLine{72.50}{22.14}{}{0}{国王。}
\BookLine{89.50}{0.00}{\heifont}{0}{阿德墨托斯：}
\BookLine{98.00}{22.14}{}{0}{欢迎您，欢迎您，尊贵的忒瑞西阿斯。您是我}
\BookLine{106.50}{22.14}{}{0}{父亲的老朋友了。欢迎来到费莱。我是塞萨利}
\BookLine{115.00}{22.14}{}{0}{国王，阿德墨托斯。}
\BookLine{132.00}{0.00}{\heifont}{0}{忒瑞西阿斯：}
\BookLine{140.50}{22.14}{}{0}{{\kaifont （舞着手杖）}后边去，往后站。又挤又乱的……}
\BookLine{149.00}{22.14}{}{0}{你，带着耳朵呢吗？}
\BookLine{166.00}{0.00}{\heifont}{0}{阿德墨托斯：}
\BookLine{174.50}{22.14}{}{0}{是的，忒瑞西阿斯。}
\newpage
% 原书页码：54
\null
\BookNumber{54}
\Vertical{4}{169}{}{6.2}{8}阿\\尔\\刻\\提\\斯\\之\\歌\EndVertical
\BookLine{30.00}{0.00}{\heifont}{0}{忒瑞西阿斯：}
\BookLine{38.50}{22.14}{}{0}{那就把你的耳屎都清理干净，仔细聆听神的}
\BookLine{47.00}{22.14}{}{0}{教诲。}
\BookLine{64.00}{0.00}{\heifont}{0}{阿德墨托斯：}
\BookLine{72.50}{22.14}{}{0}{洗耳恭听，忒瑞西阿斯。}
\BookLine{89.50}{0.00}{\heifont}{0}{忒瑞西阿斯：}
\BookLine{98.00}{22.14}{}{0}{迈锡尼国王阿特柔斯，聆听神的教诲——}
\BookLine{115.00}{0.00}{\heifont}{0}{阿德墨托斯：}
\BookLine{123.50}{22.14}{}{0}{阿特柔斯？尊贵的忒瑞西阿斯，我是阿德墨托}
\BookLine{132.00}{22.14}{}{0}{斯，是——}
\BookLine{149.00}{0.00}{\heifont}{0}{忒瑞西阿斯：}
\BookLine{157.50}{22.14}{}{0}{阿德墨托斯？哦，那阿德墨托斯，你先住嘴，}
\BookLine{166.00}{22.14}{}{0}{听我把话说完。我从德尔斐的阿波罗神殿为你}
\BookLine{174.50}{22.14}{}{0}{带来了一个消息。荣耀，一份至高无上的荣耀}
\newpage
% 原书页码：55
\null
\BookNumber{55}
\Vertical{129}{11}{}{7}{10}第\\一\\幕\EndVertical
\BookLine{30.00}{22.14}{}{0}{将会降临塞萨利。孩子，这儿是塞萨利吧？{\kaifont （把}}
\BookLine{38.50}{22.14}{}{0}{{\kaifont 手放在男孩的头上，男孩点了点头）}一份至高}
\BookLine{47.00}{22.14}{}{0}{无上的荣耀，一场危难将会降临塞萨利。}
\BookLine{64.00}{0.00}{\heifont}{0}{阿德墨托斯：}
\BookLine{72.50}{22.14}{}{0}{您说危难？忒瑞西阿斯？}
\BookLine{89.50}{0.00}{\heifont}{0}{忒瑞西阿斯：}
\BookLine{98.00}{22.14}{}{0}{荣耀，还有危难。危难就是荣耀，傻瓜。不对，}
\BookLine{106.50}{22.14}{}{0}{荣耀就是危难。在塞萨利，所有事情不都是要}
\BookLine{115.00}{22.14}{}{0}{先告诉你吗？}
\BookLine{132.00}{0.00}{\heifont}{0}{阿德墨托斯：}
\BookLine{140.50}{22.14}{}{0}{稍等片刻，忒瑞西阿斯。给我的消息同样应该}
\BookLine{149.00}{22.14}{}{0}{让我未来的王后知悉。阿格莱亚，把公主叫来。}
\BookLine{166.00}{7.38}{\kaifont}{0}{阿格莱亚匆匆走进宫殿。}
\newpage
% 原书页码：56
\null
\BookNumber{56}
\Vertical{4}{169}{}{6.2}{8}阿\\尔\\刻\\提\\斯\\之\\歌\EndVertical
\BookLine{30.00}{0.00}{\heifont}{0}{阿德墨托斯：}
\BookLine{38.50}{22.14}{}{0}{我要在今天迎娶伊俄尔科斯国王佩利阿斯的女}
\BookLine{47.00}{22.14}{}{0}{儿阿尔刻提斯，忒瑞西阿斯，您将是最尊贵的}
\BookLine{55.50}{22.14}{}{0}{客人。请先休息片刻，忒瑞西阿斯……}
\BookLine{72.50}{0.00}{\heifont}{0}{忒瑞西阿斯：}
\BookLine{81.00}{22.14}{}{0}{婚礼我天天见，让王后快点。听我说，克里特}
\BookLine{89.50}{22.14}{}{0}{国王米诺斯……孩子，他叫什么名字来着？}
\BookLine{106.50}{7.38}{\kaifont}{0}{男孩耳语。}
\BookLine{123.50}{0.00}{\heifont}{0}{忒瑞西阿斯：}
\BookLine{132.00}{22.14}{}{0}{好吧，这有什么关系？王后来了吗？}
\BookLine{149.00}{7.38}{\kaifont}{0}{阿尔刻提斯从宫中走出来，急匆匆的，并充满好奇。}
\BookLine{166.00}{0.00}{\heifont}{0}{阿德墨托斯：}
\BookLine{174.50}{22.14}{}{0}{她来了。}
\newpage
% 原书页码：57
\null
\BookNumber{57}
\Vertical{129}{11}{}{7}{10}第\\一\\幕\EndVertical
\BookLine{30.00}{0.00}{\heifont}{0}{阿尔刻提斯：}
\BookLine{38.50}{22.14}{}{0}{尊贵的忒瑞西阿斯，伟大的忒瑞西阿斯！我父}
\BookLine{47.00}{22.14}{}{0}{亲的老朋友。我是阿尔刻提斯，伊俄尔科斯国}
\BookLine{55.50}{22.14}{}{0}{王佩利阿斯的女儿。}
\BookLine{72.50}{0.00}{\heifont}{0}{忒瑞西阿斯：}
\BookLine{81.00}{22.14}{}{0}{{\kaifont （挥舞着手杖驱逐着围在他身边的人）}后边！往}
\BookLine{89.50}{22.14}{}{0}{后去！野鹅，野鸭，呱呱乱叫。闭上你们的嘴，}
\BookLine{98.00}{22.14}{}{0}{安静一会儿。众神与人类之父宙斯命令了……}
\BookLine{106.50}{22.14}{}{0}{已经命令……孩子，他命令了什么来着？}
\BookLine{123.50}{7.38}{\kaifont}{0}{男孩耳语，忒瑞西阿斯用手杖打了他一下。}
\BookLine{140.50}{0.00}{\heifont}{0}{忒瑞西阿斯：}
\BookLine{149.00}{22.14}{}{0}{嗯，你们不必……他命令阿波罗——我的主}
\BookLine{157.50}{22.14}{}{0}{人——阿波罗从奥林匹斯山下来，在人间生活}
\BookLine{166.00}{22.14}{}{0}{一年，从这个至点到下个至点，作为一个人类}
\BookLine{174.50}{22.14}{}{0}{和大家一起生活。我已经把信息带到。孩子，}
\newpage
% 原书页码：58
\null
\BookNumber{58}
\Vertical{4}{169}{}{6.2}{8}阿\\尔\\刻\\提\\斯\\之\\歌\EndVertical
\BookLine{30.00}{22.14}{}{0}{我们上路吧。{\kaifont （他转身要走）}}
\BookLine{47.00}{0.00}{\heifont}{0}{阿尔刻提斯：}
\BookLine{55.50}{22.14}{}{0}{{\kaifont （当男孩还在对忒瑞西阿斯耳语时）}阿波罗要在}
\BookLine{64.00}{22.14}{}{0}{这地上生活？}
\BookLine{81.00}{0.00}{\heifont}{0}{忒瑞西阿斯：}
\BookLine{89.50}{22.14}{}{0}{{\kaifont （对男孩）}是的，是的。我又没聋。阿波罗，我}
\BookLine{98.00}{22.14}{}{0}{的主人，他选择了生活在这里。{\kaifont （用手杖敲敲脚}}
\BookLine{106.50}{22.14}{}{0}{{\kaifont 下的土地）}在这儿，作为塞萨利国王阿德墨托}
\BookLine{115.00}{22.14}{}{0}{斯的仆人。}
\BookLine{132.00}{0.00}{\heifont}{0}{阿德墨托斯：}
\BookLine{140.50}{22.14}{}{0}{在这儿？这儿？尊贵的忒瑞西阿斯？{\kaifont （飞快走到}}
\BookLine{149.00}{22.14}{}{0}{{\kaifont 他身边）}稍等一下，忒瑞西阿斯。我该怎么理}
\BookLine{157.50}{22.14}{}{0}{解这件事情？神圣的忒瑞西阿斯，您该不是说，}
\BookLine{166.00}{22.14}{}{0}{阿波罗要在这儿做我的仆人，每一天？每一天}
\BookLine{174.50}{22.14}{}{0}{都和我们生活在一起？}
\newpage
% 原书页码：59
\null
\BookNumber{59}
\Vertical{129}{11}{}{7}{10}第\\一\\幕\EndVertical
\BookLine{30.00}{7.38}{\kaifont}{0}{大家都屏气凝神地望着他，忒瑞西阿斯却垂下手臂，像}
\BookLine{38.50}{0.00}{\kaifont}{0}{是睡着了一样。突然他醒了过来，说起话来。}
\BookLine{55.50}{0.00}{\heifont}{0}{忒瑞西阿斯：}
\BookLine{64.00}{22.14}{}{0}{门外有四个牧人。他们在这一年之中会成为你}
\BookLine{72.50}{22.14}{}{0}{的仆人。给他们找些活儿干。他们其中之一就}
\BookLine{81.00}{22.14}{}{0}{是阿波罗。}
\BookLine{98.00}{0.00}{\heifont}{0}{阿德墨托斯：}
\BookLine{106.50}{22.14}{}{0}{{\kaifont （重复道）}其中一个是阿波罗？}
\BookLine{123.50}{0.00}{\heifont}{0}{忒瑞西阿斯：}
\BookLine{132.00}{22.14}{}{0}{四个牧人。其中一个是阿波罗。不用去想哪一}
\BookLine{140.50}{22.14}{}{0}{个是神，我不知道，你也不会知道。也别再问}
\BookLine{149.00}{22.14}{}{0}{我问题了，我没什么可回答的了。孩子，把牧}
\BookLine{157.50}{22.14}{}{0}{人们叫来。}
\BookLine{174.50}{7.38}{\kaifont}{0}{男孩走出去，静场。}
\newpage
% 原书页码：60
\null
\BookNumber{60}
\Vertical{4}{169}{}{6.2}{8}阿\\尔\\刻\\提\\斯\\之\\歌\EndVertical
\BookLine{30.00}{0.00}{\heifont}{0}{阿德墨托斯：}
\BookLine{38.50}{22.14}{}{0}{忒瑞西阿斯，我们是不是应该掩面下跪？}
\BookLine{55.50}{0.00}{\heifont}{0}{忒瑞西阿斯：}
\BookLine{64.00}{22.14}{}{0}{你没有听到我刚刚说的话吗？阿波罗在这里，}
\BookLine{72.50}{22.14}{}{0}{是一个普通人类。一个普通人，和普通的牧人}
\BookLine{81.00}{22.14}{}{0}{一样。像我这样就行了！}
\BookLine{98.00}{7.38}{\kaifont}{0}{男孩回来站到忒瑞西阿斯身边。四个牧人走了进来。他}
\BookLine{106.50}{0.00}{\kaifont}{0}{们满身尘土，肮脏且不修边幅，在人群面前显得十分害羞，}
\BookLine{115.00}{0.00}{\kaifont}{0}{他们卑躬屈膝地手扶前额，转向墙角站成一排，眼睛也不知}
\BookLine{123.50}{0.00}{\kaifont}{0}{该往哪儿看。有两人带着酒囊，每人都有一支手杖。忒瑞西}
\BookLine{132.00}{0.00}{\kaifont}{0}{阿斯粗声粗气地对他们说话。}
\BookLine{149.00}{0.00}{\heifont}{0}{忒瑞西阿斯：}
\BookLine{157.50}{22.14}{}{0}{{\kaifont （继续）}过来，别磨磨蹭蹭的。给你们的新主人}
\BookLine{166.00}{22.14}{}{0}{行礼。谁都看得出来你们喝了酒，工作开了个}
\BookLine{174.50}{22.14}{}{0}{好头啊。{\kaifont （挥舞着手杖）}要是我看到了肯定会揍}
\newpage
% 原书页码：61
\null
\BookNumber{61}
\Vertical{129}{11}{}{7}{10}第\\一\\幕\EndVertical
\BookLine{30.00}{22.14}{}{0}{你们。抬脚，往前走。孩子，四个都在这儿了}
\BookLine{38.50}{22.14}{}{0}{吗？嗯，国王成哑巴啦？}
\BookLine{55.50}{0.00}{\heifont}{0}{阿德墨托斯：}
\BookLine{64.00}{22.14}{}{0}{{\kaifont （整理了下情绪）}欢迎你们到塞萨利。欢迎你们}
\BookLine{72.50}{22.14}{}{0}{来参加宴会，今天是我的婚礼。明天我就会将}
\BookLine{81.00}{22.14}{}{0}{羊群和牛群分派给你们。长途跋涉辛苦了。欢}
\BookLine{89.50}{22.14}{}{0}{迎你们到塞萨利……忒瑞西阿斯，也欢迎你，}
\BookLine{98.00}{22.14}{}{0}{长途跋涉辛苦了。需要沐浴休息一下吗？}
\BookLine{115.00}{0.00}{\heifont}{0}{忒瑞西阿斯：}
\BookLine{123.50}{22.14}{}{0}{我还要继续长途跋涉。我的信息已经带到了。}
\BookLine{132.00}{22.14}{}{0}{孩子，带路走吧。}
\BookLine{149.00}{0.00}{\heifont}{0}{阿尔刻提斯：}
\BookLine{157.50}{22.14}{}{0}{{\kaifont （几步走到他面前，低声地）}神圣的忒瑞西阿}
\BookLine{166.00}{22.14}{}{0}{斯，没有什么消息带给阿尔刻提斯吗？}
\newpage
% 原书页码：62
\null
\BookNumber{62}
\Vertical{4}{169}{}{6.2}{8}阿\\尔\\刻\\提\\斯\\之\\歌\EndVertical
\BookLine{30.00}{0.00}{\heifont}{0}{忒瑞西阿斯：}
\BookLine{38.50}{22.14}{}{0}{这个女人是谁？}
\BookLine{55.50}{0.00}{\heifont}{0}{阿尔刻提斯：}
\BookLine{64.00}{22.14}{}{0}{我是阿尔刻提斯，伊俄尔科斯国王佩利阿斯的}
\BookLine{72.50}{22.14}{}{0}{女儿。我向德尔斐送去了许多祷告和奉献。}
\BookLine{89.50}{0.00}{\heifont}{0}{忒瑞西阿斯：}
\BookLine{98.00}{22.14}{}{0}{祷告和奉献，多得数不清。孩子，带路。}
\BookLine{115.00}{7.38}{\kaifont}{0}{男孩拉住他的袖子，扶着他的肩膀，想要对他耳语几句。}
\BookLine{132.00}{0.00}{\heifont}{0}{忒瑞西阿斯：}
\BookLine{140.50}{22.14}{}{0}{哦对，我是有一个消息，要带给一个女孩，或}
\BookLine{149.00}{22.14}{}{0}{是女人。带给伊俄卡斯忒，还是阿尔刻提斯，}
\BookLine{157.50}{22.14}{}{0}{或者是得伊阿尼拉，我也记不得是谁了，不过}
\BookLine{166.00}{22.14}{}{0}{这事儿我还记着。你这孩子，别老拽我！{\kaifont （拿着}}
\BookLine{174.50}{22.14}{}{0}{{\kaifont 手杖打男孩几下）}废物！放肆！}
\newpage
% 原书页码：63
\null
\BookNumber{63}
\Vertical{129}{11}{}{7}{10}第\\一\\幕\EndVertical
\BookLine{30.00}{7.38}{\kaifont}{0}{男孩摔倒在地，忒瑞西阿斯并没有住手。}
\BookLine{47.00}{0.00}{\heifont}{0}{男孩：}
\BookLine{55.50}{22.14}{}{0}{{\kaifont （喊叫着）}忒瑞西阿斯！救命！救命！阿德墨托}
\BookLine{64.00}{22.14}{}{0}{斯国王！}
\BookLine{81.00}{0.00}{\heifont}{0}{阿德墨托斯：}
\BookLine{89.50}{22.14}{}{0}{是啊，忒瑞西阿斯，这孩子也没有——}
\BookLine{106.50}{0.00}{\heifont}{0}{阿尔刻提斯：}
\BookLine{115.00}{22.14}{}{0}{只是个小小的错误，忒瑞西阿斯。求你放过他}
\BookLine{123.50}{22.14}{}{0}{吧，他知道错了。}
\BookLine{140.50}{0.00}{\heifont}{0}{忒瑞西阿斯：}
\BookLine{149.00}{22.14}{}{0}{{\kaifont （突然停了下来，盯着阿尔刻提斯）}我不管你叫}
\BookLine{157.50}{22.14}{}{0}{什么：伊俄卡斯忒，勒达，赫耳弥俄涅——}
\newpage
% 原书页码：64
\null
\BookNumber{64}
\Vertical{4}{169}{}{6.2}{8}阿\\尔\\刻\\提\\斯\\之\\歌\EndVertical
\BookLine{30.00}{0.00}{\heifont}{0}{阿尔刻提斯：}
\BookLine{38.50}{22.14}{}{0}{阿尔刻提斯。}
\BookLine{55.50}{0.00}{\heifont}{0}{忒瑞西阿斯：}
\BookLine{64.00}{22.14}{}{0}{我是有个消息，要带给某个女孩，但我忘了，}
\BookLine{72.50}{22.14}{}{0}{要么我已经说过了。嗯，我已经说过了，借着}
\BookLine{81.00}{22.14}{}{0}{雷声与闪电，借着神圣的祭坛——如果人类学}
\BookLine{89.50}{22.14}{}{0}{不会聆听，德尔斐要来有什么用？}
\BookLine{106.50}{7.38}{\kaifont}{0}{忒瑞西阿斯跟在男孩后面走了出去，阿德墨托斯跟上}
\BookLine{115.00}{0.00}{\kaifont}{0}{几步。}
\BookLine{132.00}{0.00}{\heifont}{0}{阿德墨托斯：}
\BookLine{140.50}{22.14}{}{0}{您说，您说会有一场危难是吗，忒瑞西阿斯？}
\BookLine{157.50}{0.00}{\heifont}{0}{忒瑞西阿斯：}
\BookLine{166.00}{22.14}{}{0}{{\kaifont （马上要走出门外）}当然有危难了，呆子。每次}
\BookLine{174.50}{22.14}{}{0}{他们{\kaifont （指着天空）}的来临就是一场危难。}
\newpage
% 原书页码：65
\null
\BookNumber{65}
\Vertical{129}{11}{}{7}{10}第\\一\\幕\EndVertical
\BookLine{30.00}{0.00}{\heifont}{0}{阿德墨托斯：}
\BookLine{38.50}{22.14}{}{0}{可我父亲说，阿波罗是爱着塞萨利的……}
\BookLine{55.50}{0.00}{\heifont}{0}{忒瑞西阿斯：}
\BookLine{64.00}{22.14}{}{0}{是的，爱爱爱。大家还是自个爱自个吧。你看}
\BookLine{72.50}{22.14}{}{0}{看我，五六百岁了，神很爱我，死都不让我死。}
\BookLine{81.00}{22.14}{}{0}{如果神不爱人类，我们肯定都很快乐；换个角}
\BookLine{89.50}{22.14}{}{0}{度同样成立：如果我们不爱神，我们肯定也很}
\BookLine{98.00}{22.14}{}{0}{快乐。}
\BookLine{115.00}{7.38}{\kaifont}{0}{忒瑞西阿斯带着男孩下场。一个尴尬的停顿。}
\BookLine{123.50}{7.38}{\kaifont}{0}{阿德墨托斯整理了一下情绪，带着实事求是却又威严的}
\BookLine{132.00}{0.00}{\kaifont}{0}{语气对四个牧人开始讲话。}
\BookLine{149.00}{0.00}{\heifont}{0}{阿德墨托斯：}
\BookLine{157.50}{22.14}{}{0}{再次欢迎你们来到塞萨利，来到费莱。{\kaifont （对看门}}
\BookLine{166.00}{22.14}{}{0}{{\kaifont 人）}请确保安顿好他们。{\kaifont （对牧人们）}很高兴你}
\BookLine{174.50}{22.14}{}{0}{们也成为了今天婚礼的客人。}
\newpage
% 原书页码：66
\null
\BookNumber{66}
\Vertical{4}{169}{}{6.2}{8}阿\\尔\\刻\\提\\斯\\之\\歌\EndVertical
\BookLine{30.00}{7.38}{\kaifont}{0}{看门人带领四个牧人沿着小路走向泉边，阿德墨托斯和}
\BookLine{38.50}{0.00}{\kaifont}{0}{舞台上的人们都好奇地看着。他们走过阿德墨托斯身边时，}
\BookLine{47.00}{0.00}{\kaifont}{0}{都像仆人一样怯生生的，到了泉边他们才放松下来，互相传}
\BookLine{55.50}{0.00}{\kaifont}{0}{递着酒囊，其中一人渐渐睡着了。阿德墨托斯并没有看向阿}
\BookLine{64.00}{0.00}{\kaifont}{0}{尔刻提斯，他对她说话，也是在对自己说。}
\BookLine{81.00}{0.00}{\heifont}{0}{阿德墨托斯：}
\BookLine{89.50}{22.14}{}{0}{{\kaifont （继续）}我也不知道该怎么理解这些事儿……我}
\BookLine{98.00}{22.14}{}{0}{自己也不过就是个牧人。阿尔刻提斯，费莱很}
\BookLine{106.50}{22.14}{}{0}{需要你。}
\BookLine{123.50}{7.38}{\kaifont}{0}{他伸出双手背到身后。她，沉思着愣在原地，毫无回应。}
\BookLine{140.50}{0.00}{\heifont}{0}{阿德墨托斯：}
\BookLine{149.00}{22.14}{}{0}{我得回去招呼客人了。{\kaifont （再次表达畏惧）}我真的}
\BookLine{157.50}{22.14}{}{0}{不知道，该怎么理解这些事情。{\kaifont （忽然笑了一}}
\BookLine{166.00}{22.14}{}{0}{{\kaifont 下）}阿尔刻提斯，在塞萨利有一个古老的习俗，}
\BookLine{174.50}{22.14}{}{0}{直到婚礼当晚，新郎都不能见到新娘的面容。}
\newpage
% 原书页码：67
\null
\BookNumber{67}
\Vertical{129}{11}{}{7}{10}第\\一\\幕\EndVertical
\BookLine{30.00}{22.14}{}{0}{这个习俗已经流传好几百年了，在伊俄尔科斯}
\BookLine{38.50}{22.14}{}{0}{也是这样吗？}
\BookLine{55.50}{0.00}{\heifont}{0}{阿尔刻提斯：}
\BookLine{64.00}{22.14}{}{0}{{\kaifont （低声）}是的，阿德墨托斯。}
\BookLine{81.00}{0.00}{\heifont}{0}{阿德墨托斯：}
\BookLine{89.50}{22.14}{}{0}{{\kaifont （活力满满地走过她身边，用手挡着自己的脸）}}
\BookLine{98.00}{22.14}{}{0}{很快——感谢神的恩赐——我就可以一直望着}
\BookLine{106.50}{22.14}{}{0}{你，直到我死去。}
\BookLine{123.50}{7.38}{\kaifont}{0}{走到王宫门口时，他忽然想起了什么。他仍然用手挡着}
\BookLine{132.00}{0.00}{\kaifont}{0}{自己的脸，缓缓地走到四个牧人在泉边所站立的地方。他下}
\BookLine{140.50}{0.00}{\kaifont}{0}{定决心般深深呼吸一口气，平静而直率地说道。}
\BookLine{157.50}{0.00}{\heifont}{0}{阿德墨托斯：}
\BookLine{166.00}{22.14}{}{0}{阿波罗，我父亲、我祖先、我土地的朋友，我}
\BookLine{174.50}{22.14}{}{0}{是一个虔诚的普通人，但我对真正的虔诚知晓}
\newpage
% 原书页码：68
\null
\BookNumber{68}
\Vertical{4}{169}{}{6.2}{8}阿\\尔\\刻\\提\\斯\\之\\歌\EndVertical
\BookLine{30.00}{22.14}{}{0}{不多。如果我不经意间有所冒犯，或是犯了错}
\BookLine{38.50}{22.14}{}{0}{误，希望这位“祂”，我的家族、我的人民的}
\BookLine{47.00}{22.14}{}{0}{朋友，能够原谅我。您在我人生最幸福的一}
\BookLine{55.50}{22.14}{}{0}{天到来，也请继续赐福于我和我的后裔吧……}
\BookLine{64.00}{22.14}{}{0}{{\kaifont （停顿）}我不太会说话，不过您能看透人心，请}
\BookLine{72.50}{22.14}{}{0}{鉴察我的心，或者，您可以在我脑中埋下愿望，}
\BookLine{81.00}{22.14}{}{0}{那些只有您能满足的愿望。}
\BookLine{98.00}{7.38}{\kaifont}{0}{他转过身迅速地走进宫殿。}
\BookLine{115.00}{0.00}{\heifont}{0}{阿尔刻提斯：}
\BookLine{123.50}{22.14}{}{0}{{\kaifont （一直在盯着牧人们看，带着期盼和渴望，又带}}
\BookLine{132.00}{22.14}{}{0}{{\kaifont 着怀疑和抗拒——尽管此刻他们已经从她身前}}
\BookLine{140.50}{22.14}{}{0}{{\kaifont 躲开了）}阿波罗在这儿吗？}
\BookLine{157.50}{0.00}{\heifont}{0}{阿格莱亚：}
\BookLine{166.00}{22.14}{}{0}{公主！}
\newpage
% 原书页码：69
\null
\BookNumber{69}
\Vertical{129}{11}{}{7}{10}第\\一\\幕\EndVertical
\BookLine{30.00}{0.00}{\heifont}{0}{阿尔刻提斯：}
\BookLine{38.50}{22.14}{}{0}{就是他们其中之一？那个老头就是德尔斐的}
\BookLine{47.00}{22.14}{}{0}{忒瑞西阿斯——那个颤颤巍巍、疯疯癫癫的}
\BookLine{55.50}{22.14}{}{0}{老头？}
\BookLine{72.50}{0.00}{\heifont}{0}{阿格莱亚：}
\BookLine{81.00}{22.14}{}{0}{{\kaifont （被惊讶到，坚定地）}不要怀疑这些事情，}
\BookLine{89.50}{22.14}{}{0}{公主。}
\BookLine{106.50}{0.00}{\heifont}{0}{阿尔刻提斯：}
\BookLine{115.00}{22.14}{}{0}{{\kaifont （朝着牧人们走了几步，突然做了决定）}让我和}
\BookLine{123.50}{22.14}{}{0}{他们单独待会儿。}
\BookLine{140.50}{7.38}{\kaifont}{0}{阿格莱亚对看门人做了个手势，两人下场。在接下来的}
\BookLine{149.00}{0.00}{\kaifont}{0}{谈话中，一个牧人睡着了，还在打鼾，其他几个则在她面前}
\BookLine{157.50}{0.00}{\kaifont}{0}{十分窘迫。他们互相递着酒囊，都不敢抬头看她。}
\newpage
% 原书页码：70
\null
\BookNumber{70}
\Vertical{4}{169}{}{6.2}{8}阿\\尔\\刻\\提\\斯\\之\\歌\EndVertical
\BookLine{30.00}{0.00}{\heifont}{0}{阿尔刻提斯：}
\BookLine{38.50}{22.14}{}{0}{您在这儿吗？我已经向您呼求成百上千次}
\BookLine{47.00}{22.14}{}{0}{了——向着天空、群星、太阳。我向德尔斐献}
\BookLine{55.50}{22.14}{}{0}{上了数不清的祷告。您现在是否，确确实实，}
\BookLine{64.00}{22.14}{}{0}{听到了我的声音？{\kaifont （静场，突然被鼾声打断）}有}
\BookLine{72.50}{22.14}{}{0}{人说，您不存在。有人说，神明都在很远的地}
\BookLine{81.00}{22.14}{}{0}{方，他们在奥林匹斯山上欢宴，饮酒，沉睡。}
\BookLine{89.50}{22.14}{}{0}{我已经把生命献给了您。您知道，我希望只为}
\BookLine{98.00}{22.14}{}{0}{您而活：去学习——由您来教导——生命的意}
\BookLine{106.50}{22.14}{}{0}{义。{\kaifont （没有回应）}人类是不是无法得到启示、教}
\BookLine{115.00}{22.14}{}{0}{训？我们是不是已经被抛弃了？{\kaifont （她等待着，场}}
\BookLine{123.50}{22.14}{}{0}{{\kaifont 上只有牧人们尴尬的沉默。她转身向王宫走去，}}
\BookLine{132.00}{22.14}{}{0}{{\kaifont 痛苦地自言自语）}那么我们只能自己寻找出路}
\BookLine{140.50}{22.14}{}{0}{了……生命不过是追寻这些那些过后的虚无，}
\BookLine{149.00}{22.14}{}{0}{是一种充满激情的空虚……}
\BookLine{166.00}{7.38}{\kaifont}{0}{第一个牧人——最脏的，四人中最不起眼的——站了起}
\BookLine{174.50}{0.00}{\kaifont}{0}{来。他谦卑而害羞地摸了摸自己的帽檐，说道。}
\newpage
% 原书页码：71
\null
\BookNumber{71}
\Vertical{129}{11}{}{7}{10}第\\一\\幕\EndVertical
\BookLine{30.00}{0.00}{\heifont}{0}{牧人：}
\BookLine{38.50}{22.14}{}{0}{公主，那个老头说我们之中有一个是神？他确}
\BookLine{47.00}{22.14}{}{0}{实这么说了？阿波罗？如果这样，夫人，我和}
\BookLine{55.50}{22.14}{}{0}{您一样惊讶。夫人，我们四人在过去三十天穿}
\BookLine{64.00}{22.14}{}{0}{越整个希腊，用同一个酒囊饮酒，用同一个碟}
\BookLine{72.50}{22.14}{}{0}{子吃饭，围着同一个火堆睡觉。如果我们之中}
\BookLine{81.00}{22.14}{}{0}{真有一位神明，我怎么会没发现呢？}
\BookLine{98.00}{7.38}{\kaifont}{0}{阿尔刻提斯带着希望和厌恶，朝他走了几步。}
\BookLine{115.00}{0.00}{\heifont}{0}{牧人：}
\BookLine{123.50}{22.14}{}{0}{照我看，夫人，我发誓我们只是普通的牧人罢}
\BookLine{132.00}{22.14}{}{0}{了，无知的牧人。但……但是有件事值得一}
\BookLine{140.50}{22.14}{}{0}{提，夫人，我们也不是平庸的牧人。怎么说}
\BookLine{149.00}{22.14}{}{0}{呢？那边那个家伙——在打鼾的那个家伙，没}
\BookLine{157.50}{22.14}{}{0}{有什么疾病是他治不好的，被蛇咬伤或是背部}
\BookLine{166.00}{22.14}{}{0}{骨折，不过我知道他不是神，公主。他旁边的}
\BookLine{174.50}{22.14}{}{0}{那个家伙，那个！{\kaifont （他走过去踢了那个牧人一}}
\newpage
% 原书页码：72
\null
\BookNumber{72}
\Vertical{4}{169}{}{6.2}{8}阿\\尔\\刻\\提\\斯\\之\\歌\EndVertical
\BookLine{30.00}{22.14}{}{0}{{\kaifont 脚）}你能不能别喝了？公主看着你呢。他从来}
\BookLine{38.50}{22.14}{}{0}{不会迷路，即使在最黑的夜里，他也能准确分}
\BookLine{47.00}{22.14}{}{0}{辨东西南北，很神奇。不过我知道他也不是太}
\BookLine{55.50}{22.14}{}{0}{阳之神。{\kaifont （小声说道）}另外，他有些很下流的习}
\BookLine{64.00}{22.14}{}{0}{惯，很下流。}
\BookLine{81.00}{0.00}{\heifont}{0}{阿尔刻提斯：}
\BookLine{89.50}{22.14}{}{0}{{\kaifont （毫不在意）}那个呢？}
\BookLine{106.50}{0.00}{\heifont}{0}{牧人：}
\BookLine{115.00}{22.14}{}{0}{那个？他是我们的歌手。}
\BookLine{132.00}{0.00}{\heifont}{0}{阿尔刻提斯：}
\BookLine{140.50}{22.14}{}{0}{啊！}
\BookLine{157.50}{0.00}{\heifont}{0}{牧人：}
\BookLine{166.00}{22.14}{}{0}{我跟您说，当他弹着竖琴唱起歌时——哦，公}
\BookLine{174.50}{22.14}{}{0}{主！每当那时我都对自己说“这是神啊”。即便}
\newpage
% 原书页码：73
\null
\BookNumber{73}
\Vertical{129}{11}{}{7}{10}第\\一\\幕\EndVertical
\BookLine{30.00}{22.14}{}{0}{没有什么值得快乐或忧伤的事情发生，他的音}
\BookLine{38.50}{22.14}{}{0}{乐也会让你感到满溢的欢喜和忧愁。他可以使}
\BookLine{47.00}{22.14}{}{0}{爱的记忆变得比爱本身更加甜蜜，但是公主，}
\BookLine{55.50}{22.14}{}{0}{他不是神。{\kaifont （她早就想反驳他，想要争辩一番）}}
\BookLine{64.00}{22.14}{}{0}{一个总是一副苦脸，随时要把自己喝死的人怎}
\BookLine{72.50}{22.14}{}{0}{么会是神呢？他在慢性自杀，您知道的，在我}
\BookLine{81.00}{22.14}{}{0}{们眼皮底下。神明们可不憎恨自己，公主。}
\BookLine{98.00}{0.00}{\heifont}{0}{阿尔刻提斯：}
\BookLine{106.50}{22.14}{}{0}{你呢？}
\BookLine{123.50}{0.00}{\heifont}{0}{牧人：}
\BookLine{132.00}{22.14}{}{0}{我？我是阿波罗？我不但不是阿波罗，我也不}
\BookLine{140.50}{22.14}{}{0}{相信阿波罗在这儿。}
\BookLine{157.50}{0.00}{\heifont}{0}{阿尔刻提斯：}
\BookLine{166.00}{22.14}{}{0}{忒瑞西阿斯，忒瑞西阿斯说了……}
\newpage
% 原书页码：74
\null
\BookNumber{74}
\Vertical{4}{169}{}{6.2}{8}阿\\尔\\刻\\提\\斯\\之\\歌\EndVertical
\BookLine{30.00}{0.00}{\heifont}{0}{牧人：}
\BookLine{38.50}{22.14}{}{0}{那个忒瑞西阿斯——那个路都走不稳还稀里糊}
\BookLine{47.00}{22.14}{}{0}{涂的老头？还能找到比他更差劲的送信人吗？}
\BookLine{55.50}{22.14}{}{0}{还能把神谕说得再不清不楚一点吗？}
\BookLine{72.50}{7.38}{\kaifont}{0}{阿尔刻提斯转身向王宫走去。她又转过身来看着牧人，}
\BookLine{81.00}{0.00}{\kaifont}{0}{像是在对自己说话。}
\BookLine{98.00}{0.00}{\heifont}{0}{阿尔刻提斯：}
\BookLine{106.50}{22.14}{}{0}{这样说来，我们真的很可悲。我们不仅得不到}
\BookLine{115.00}{22.14}{}{0}{帮助，还傻傻地相信我们可能会得到一些帮助。}
\BookLine{132.00}{0.00}{\heifont}{0}{牧人：}
\BookLine{140.50}{22.14}{}{0}{{\kaifont （向前走了几步到舞台上）}但如果神明们真的存}
\BookLine{149.00}{22.14}{}{0}{在，他们会如何对我们说话呢？他们会使用哪}
\BookLine{157.50}{22.14}{}{0}{种语言？和他们比起来，我们不过是病态的、}
\BookLine{166.00}{22.14}{}{0}{将死的、又聋又瞎还忙得团团转的小丑。为什}
\BookLine{174.50}{22.14}{}{0}{么有人会觉得神爱我们。您想得通吗？当爱人}
\newpage
% 原书页码：75
\null
\BookNumber{75}
\Vertical{129}{11}{}{7}{10}第\\一\\幕\EndVertical
\BookLine{30.00}{22.14}{}{0}{们之间隔着如此巨大的鸿沟，公主，这算哪门}
\BookLine{38.50}{22.14}{}{0}{子爱啊？{\kaifont （他慢慢走回自己的位置）}这种爱不会}
\BookLine{47.00}{22.14}{}{0}{快乐，简直不用怀疑。}
\BookLine{64.00}{0.00}{\heifont}{0}{阿尔刻提斯：}
\BookLine{72.50}{22.14}{}{0}{{\kaifont （恳切而尖锐地）}不，谁说不会快乐！}
\BookLine{89.50}{0.00}{\heifont}{0}{牧人：}
\BookLine{98.00}{22.14}{}{0}{{\kaifont （带着同样的情绪）}是。如果他真让我们看到他}
\BookLine{106.50}{22.14}{}{0}{的荣光，那会杀死我们的。{\kaifont （停顿）}对于今早发}
\BookLine{115.00}{22.14}{}{0}{生的事，我有些想法，或许有别的方法，去填}
\BookLine{123.50}{22.14}{}{0}{补这个鸿沟。或许神明们会找到一种方法，让}
\BookLine{132.00}{22.14}{}{0}{他们所爱之人变得更好，更接近他们。如果忒}
\BookLine{140.50}{22.14}{}{0}{瑞西阿斯说得没错，那阿波罗确实在塞萨利。}
\BookLine{149.00}{22.14}{}{0}{也或许，那个老糊涂捎错了信儿，或许信息本}
\BookLine{157.50}{22.14}{}{0}{来是，阿波罗在这里，分散在了每一个人身}
\BookLine{166.00}{22.14}{}{0}{上——我们四个牧人，还有其他人！比如说阿}
\BookLine{174.50}{22.14}{}{0}{德墨托斯，我才见过他几小时，我必须要说，}
\newpage
% 原书页码：76
\null
\BookNumber{76}
\Vertical{4}{169}{}{6.2}{8}阿\\尔\\刻\\提\\斯\\之\\歌\EndVertical
\BookLine{30.00}{22.14}{}{0}{公主，一开始，我对阿德墨托斯非常失望，他}
\BookLine{38.50}{22.14}{}{0}{看上去一点也没有什么不凡之处。您见过赫拉}
\BookLine{47.00}{22.14}{}{0}{克勒斯吧，公主？}
\BookLine{64.00}{0.00}{\heifont}{0}{阿尔刻提斯：}
\BookLine{72.50}{22.14}{}{0}{{\kaifont （轻轻点了点头）}见过。}
\BookLine{89.50}{0.00}{\heifont}{0}{牧人：}
\BookLine{98.00}{22.14}{}{0}{{\kaifont （突然回想起来）}啊，对。他追求过您。很有个}
\BookLine{106.50}{22.14}{}{0}{样儿！赫拉克勒斯，宙斯与阿尔克墨涅之子，您}
\BookLine{115.00}{22.14}{}{0}{一眼就能看出来。我见过不少比阿德墨托斯更强}
\BookLine{123.50}{22.14}{}{0}{的人，但渐渐地，我发现阿德墨托斯身上有一些}
\BookLine{132.00}{22.14}{}{0}{别的英雄没有的东西……这个世界在改变，慢}
\BookLine{140.50}{22.14}{}{0}{慢地改变。除非有不一样的男人、女人们出现，}
\BookLine{149.00}{22.14}{}{0}{公主，这个世界怎么会变得更好呢？}
\BookLine{166.00}{7.38}{\kaifont}{0}{阿德墨托斯上场，右肩上披着像阿波罗一样的浅蓝色的}
\BookLine{174.50}{0.00}{\kaifont}{0}{斗篷，他站在最上面一级台阶观望着。}
\newpage
% 原书页码：77
\null
\BookNumber{77}
\Vertical{129}{11}{}{7}{10}第\\一\\幕\EndVertical
\BookLine{30.00}{0.00}{\heifont}{0}{牧人：}
\BookLine{38.50}{22.14}{}{0}{这不就是，或许，就是那些伤感的爱人们，{\kaifont （指}}
\BookLine{47.00}{22.14}{}{0}{{\kaifont 指天空）}在想方设法，在我刚说的鸿沟上面搭}
\BookLine{55.50}{22.14}{}{0}{建的一座桥梁？{\kaifont （意识到阿德墨托斯的出现，谦}}
\BookLine{64.00}{22.14}{}{0}{{\kaifont 卑地低下头，用手触摸着前额，默念着）}我们}
\BookLine{72.50}{22.14}{}{0}{祝您幸福，长长久久。{\kaifont （回到自己的位置上）}}
\BookLine{89.50}{7.38}{\kaifont}{0}{太阳渐渐落山，宫殿背后的天空被染上了色彩。}
\BookLine{106.50}{0.00}{\heifont}{0}{阿德墨托斯：}
\BookLine{115.00}{22.14}{}{0}{{\kaifont （板着脸）}阿尔刻提斯公主，阿格莱亚已经告知}
\BookLine{123.50}{22.14}{}{0}{了我您离开的意愿。您，当然有离开的自由，}
\BookLine{132.00}{22.14}{}{0}{我们不会强留。塞萨利没有奴隶，公主，它的}
\BookLine{140.50}{22.14}{}{0}{王后也不会成为奴隶。我已经命令您的马夫和}
\BookLine{149.00}{22.14}{}{0}{女仆为您打点好旅途上需要的一切。{\kaifont （向下走一}}
\BookLine{157.50}{22.14}{}{0}{{\kaifont 步）}在您离开之前，我想说一件事，说这些不}
\BookLine{166.00}{22.14}{}{0}{是为了让您可怜我，也没有劝您回头的打算。}
\BookLine{174.50}{22.14}{}{0}{我说这些话是因为我们都是成年人了，对心中}
\newpage
% 原书页码：78
\null
\BookNumber{78}
\Vertical{4}{169}{}{6.2}{8}阿\\尔\\刻\\提\\斯\\之\\歌\EndVertical
\BookLine{30.00}{22.14}{}{0}{所想的事情不必掩藏。{\kaifont （轻微的停顿）}能够同时}
\BookLine{38.50}{22.14}{}{0}{驾驭狮子和野猪，这件事情直到现在都让我觉}
\BookLine{47.00}{22.14}{}{0}{得不可思议。但为什么能做到这一点，原因我}
\BookLine{55.50}{22.14}{}{0}{毫不怀疑：我爱您。我今后永远都见不到您了，}
\BookLine{64.00}{22.14}{}{0}{我永远不会再娶妻，我也永远不会是现在的我。}
\BookLine{72.50}{22.14}{}{0}{从今天开始，我不得不意识到，这个世界总是}
\BookLine{81.00}{22.14}{}{0}{有所残缺的。我是一个普通的人，但那充满在}
\BookLine{89.50}{22.14}{}{0}{我心中的爱，不是普通的爱。如果这样的爱都}
\BookLine{98.00}{22.14}{}{0}{得不到回应，那生命本身就成了一场谎言。人}
\BookLine{106.50}{22.14}{}{0}{们只需要随波逐流地生活着，所谓正义、荣誉、}
\BookLine{115.00}{22.14}{}{0}{爱，不过是短暂的、瞬间的一种妄想。愿您旅}
\BookLine{123.50}{22.14}{}{0}{途一切平安，阿尔刻提斯公主。}
\BookLine{140.50}{7.38}{\kaifont}{0}{阿尔刻提斯站在那里，慢慢抬起头，但仍不抬眼。她伸}
\BookLine{149.00}{0.00}{\kaifont}{0}{出了她的手。}
\BookLine{166.00}{0.00}{\heifont}{0}{阿尔刻提斯：}
\BookLine{174.50}{22.14}{}{0}{阿德墨托斯。{\kaifont （轻微停顿）}阿德墨托斯，再向我}
\newpage
% 原书页码：79
\null
\BookNumber{79}
\Vertical{129}{11}{}{7}{10}第\\一\\幕\EndVertical
\BookLine{30.00}{22.14}{}{0}{求一次婚吧。}
\BookLine{47.00}{7.38}{\kaifont}{0}{阿德墨托斯一个箭步向阿尔刻提斯走去，她再次伸出手，}
\BookLine{55.50}{0.00}{\kaifont}{0}{让他停下。}
\BookLine{72.50}{22.14}{}{0}{让我去爱你所爱，成为你的塞萨利的王后。让}
\BookLine{81.00}{22.14}{}{0}{我为你生儿育女，即使经历痛苦，让我在盛大}
\BookLine{89.50}{22.14}{}{0}{的宴会上挽着你走过，让我在你失意时安慰你，}
\BookLine{98.00}{22.14}{}{0}{让我在你远行归来时准备好沐浴的热水，让我}
\BookLine{106.50}{22.14}{}{0}{好好照料你的家，阿德墨托斯，像你所期待}
\BookLine{115.00}{22.14}{}{0}{的一样。让我为你、你的孩子、你的人民而生}
\BookLine{123.50}{22.14}{}{0}{活——让我每一刻都为你活着，每一刻都可以}
\BookLine{132.00}{22.14}{}{0}{为你死去。}
\BookLine{149.00}{0.00}{\heifont}{0}{阿德墨托斯：}
\BookLine{157.50}{22.14}{}{0}{{\kaifont （开心地，大声地）}不，不，阿尔刻提斯！应该}
\BookLine{166.00}{22.14}{}{0}{是我……}
\newpage
% 原书页码：80
\null
\BookNumber{80}
\Vertical{4}{169}{}{6.2}{8}阿\\尔\\刻\\提\\斯\\之\\歌\EndVertical
\BookLine{30.00}{7.38}{\kaifont}{0}{他拉起她伸出的手，他们并没有拥抱。}
\BookLine{47.00}{22.14}{}{0}{阿尔刻提斯，你是否愿意成为塞萨利国王——}
\BookLine{55.50}{22.14}{}{0}{阿德墨托斯的妻子？}
\BookLine{72.50}{0.00}{\heifont}{0}{阿尔刻提斯：}
\BookLine{81.00}{22.14}{}{0}{我愿交出我的所有，阿德墨托斯……}
\BookLine{98.00}{7.38}{\kaifont}{0}{他们走进宫殿。}
\end{document}
